% concatenated from writeup.tex
%-----------------------------------------------------------------------
% Beginning of tran-l-template.tex
%-----------------------------------------------------------------------
%
%     This is a topmatter template file for TRAN for use with AMS-LaTeX.
%
%     Templates for various common text, math and figure elements are
%     given following the \end{document} line.
%
%%%%%%%%%%%%%%%%%%%%%%%%%%%%%%%%%%%%%%%%%%%%%%%%%%%%%%%%%%%%%%%%%%%%%%%%

%     Remove any commented or uncommented macros you do not use.

\documentclass[11pt]{tran-l}
\usepackage[a4paper, top=1.5in, bottom=1.5in, left=1in, right=1in]{geometry}

%    If you need symbols beyond the basic set, uncomment this command.
\usepackage{amssymb}

%     If your article includes graphics, uncomment this command.
%\usepackage{graphicx}

%     If the article includes commutative diagrams, ...
%\usepackage[cmtip,all]{xy}


%     Update the information and uncomment if AMS is not the copyright
%     holder.
%\copyrightinfo{2009}{American Mathematical Society}

\newtheorem{theorem}{Theorem}[section]
\newtheorem{theoremL}{Theorem}[section]    % lettered version
\renewcommand{\thetheoremL}{\Alph{theoremL}}
\newtheorem{lemma}[theorem]{Lemma}
\newtheorem{prop}[theorem]{Proposition}
\newtheorem{cor}[theorem]{Corollary}

\theoremstyle{definition}
\newtheorem{definition}[theorem]{Definition}
\newtheorem{example}[theorem]{Example}
\newtheorem{xca}[theorem]{Exercise}

\theoremstyle{remark}
\newtheorem{remark}[theorem]{Remark}


\numberwithin{equation}{section}

\begin{document}



% \title[short text for running head]{full title}
\title[Large fluctuations of random multiplicative functions]{
%Almost sure 
Large fluctuations of 
%some 
sums of a random multiplicative function}


\author{Besfort Shala}
\address{Fry Building, School of Mathematics, Woodland Road, Bristol BS8 1UG, United Kingdom}
\curraddr{}
\email{besfort.shala@bristol.ac.uk}
\thanks{}



%    \subjclass is required.
%\subjclass[2010]{Primary }

\date{}

\dedicatory{}

%    Abstract is required.
\begin{abstract} Let $f$ be a Rademacher or Steinhaus random multiplicative function. For various arithmetically interesting subsets $\mathcal A\subseteq [1, N]\cap\mathbb N$ such that the distribution of $\sum_{n\in \mathcal A} f(n)$ is approximately Gaussian, we develop a general framework to understand the large fluctuations of the sum. This extends the general central limit theorem framework of Soundararajan and Xu.

In the case when $\mathcal A = (N-H, N]$ is a short interval with admissible $H=H(N)$, we show that almost surely
\begin{equation*}
    \limsup_{N\to\infty} \frac{\big\lvert\sum_{N-H<n\leq N} f(n)\big\rvert}{\sqrt{H\log \frac{N}H{}}}>0.
\end{equation*} When $\mathcal A$ is the set of values of an admissible polynomial $P\in\mathbb Z[x]$, we extend work of Klurman, Shkredov, and Xu, as well as Chinis and the author, showing that almost surely 
\begin{equation*}
   \limsup_{N\to\infty} \frac{\big\lvert\sum_{n\leq N} f(P(n))\big\rvert}{\sqrt{N \log\log N}}>0,
\end{equation*} even when $P$ is a product of linear factors over $\mathbb Q$. In this case, we also establish the corresponding almost sure upper bound, matching the law of iterated logarithm. 

An important ingredient in our work is bounding the Kantorovich--Wasserstein distance by means of a quantitative martingale central limit theorem. 
\end{abstract}

\maketitle

\section{Introduction} 

Since their introduction by Wintner \cite{Wintner1944}, random multiplicative functions have attracted a lot of attention in number theory. 
\begin{definition}
    A \emph{Steinhaus} random multiplicative function $f$ is a sequence of random variables $f(1), f(2), \ldots$ such that $f(mn)=f(m)f(n)$ for all positive integers $m$ and $n$, and for each prime $p$, $f(p)$ is uniformly distributed on the unit circle.

    A \emph{Rademacher} random multiplicative function $f$ is a sequence of random variables $f(1), f(2),$ $\ldots$  supported on square-free integers such that $f(mn) = f(m)f(n)$ whenever $m$ and $n$ are coprime, and for each prime $p$, $f(p)$ is uniformly distributed on $\{-1, 1\}$. 
\end{definition}
Although random multiplicative functions are useful in modeling deterministic objects such as Dirichlet and Archimedean characters, they have become increasingly interesting to study in their own right. Building on earlier work of Hal\'asz \cite{Halasz1977} and Hal\'asz and R\'enyi \cite{HalaszRenyi1961}, partial sums of a random multiplicative function $\sum_{n\leq N} f(n)$ have been extensively studied by Harper. Being a sum of random variables, it is natural to ask whether $\frac{1}{\sqrt N}\sum_{n\leq N} f(n)$ has an approximately Gaussian distribution as $N\to\infty$. 
Harper \cite{Harper2013} showed that this is not the case and in fact, he later showed \cite{Harper2020_MomentsRandomMultiplicativeI} that the typical size of these sums is of order $\sqrt N/(\log\log N)^{\frac14}$. One might therefore hope to obtain 
an interesting limiting distribution after normalizing by this factor instead. This is a delicate matter closely related to Gaussian multiplicative chaos, and has been resolved recently in the work of Gorodetsky and Wong \cite{gorodetsky_wong_2025_limiting} in the Steinhaus case. 

From an arithmetic point of view, it is interesting to understand sums of a random multiplicative function when the summands are restricted to subsets of positive integers. This will be our focus in this paper. Examples of such restrictions that have been considered previously include: 
\begin{itemize}\item[(i)] integers with an atypically small number of prime factors (Hough \cite{Hough2011}, Harper \cite{Harper2013}), 
\item[(ii)] integers in a short interval (Chatterjee and Soundararajan \cite{ChatterjeeSoundararajan2012}, Soundararajan and Xu \cite{soundararajan2023clt}, Pandey, Wang, and Xu \cite{PandeyWangXu2024}), 
\item[(iii)] sums of two squares in a short interval and shifted primes $p-k$ for $k\neq 0$ (Soundararajan and Xu \cite{soundararajan2023clt}), 
\item[(iv)] polynomial images (Najnudel \cite{najnudel2020consecutive}, Klurman, Shkredov, and Xu \cite{KlurmanShkredovXu2023}, Chinis and the author \cite{ChinisShala2025}).
\end{itemize} 
In each of the aforementioned works, the distribution of the sums turns out to be approximately Gaussian, demonstrating that multiplicativity seems to interfere less with independence when the summands are restricted to a less ``multiplicatively structured'' set. In this paper, we extend the general framework of Soundararajan and Xu \cite{soundararajan2023clt} on central limit theorems for random multiplicative functions to the setting of large fluctuations.

There has been considerable work in understanding almost sure large fluctuations for the full sum $\sum_{n\leq N} f(n)$ and related variants. Harper \cite{Harper2023_ASLF} showed that for any (slowly growing) function $V(N)\to\infty$, almost surely there is a sequence of positive integers $N_k\to\infty$ such that 
\begin{equation*}
    \left\lvert\sum_{n\leq N_k}f(n)\right\rvert\gg\sqrt{N_k}\frac{(\log\log N_k)^{\frac14}}{V(N_k)}.
\end{equation*}
This matches what one expects from the law of iterated logarithm, up to the factor $V(N)$. Proving the corresponding upper bound remains a challenge. Partial results with an upper bound of the right shape (namely $\sqrt N$ times a power of $\log\log N$) have been proved by Lau, Tenenbaum and Wu \cite{LauTenenbaumWu2013MeanValuesRMF}, as well as most recently by Caich \cite{Caich2023} with the current best known almost sure bound
\begin{equation*}
    \left\lvert\sum_{n\leq N} f(n)\right\rvert\ll_{\epsilon} \sqrt N(\log\log N)^{\frac{3}{4}+\epsilon}.
\end{equation*}
In the case when the summands are demanded to have a prime factor bigger than $\sqrt N$, Mastrostefano \cite{Mastrostefano2022} proved an essentially sharp almost sure upper bound with the exponent $\frac{1}{4}+\epsilon$ in place of $\frac{3}4$, with the corresponding almost sure lower bound following by work of Harper \cite{Harper2023_ASLF}; see also the work of Hardy \cite{Hardy2025LargePrime} on the distributional aspect in this case. A weighted variant $\sum_{n\leq N} \frac{f(n)}{\sqrt n}$
has been considered by Aymone, Heap and Zhao \cite{aymone2021partial}, and Hardy \cite{Hardy2023AlmostSure} in the Steinhaus case who obtained essentially sharp bounds, and Atherfold \cite{Atherfold2025} in the Rademacher case albeit without sharp bounds.

Recently, Hoban, Shah, Ismail, Verreault, and Zaman \cite{HobanShahIsmailVerreaultZaman2025} have extended the framework of Soundararajan and Xu \cite{soundararajan2023clt} in a different direction, namely in the function field setting. It is likely that the methods of this paper, combined with the work in \cite{HobanShahIsmailVerreaultZaman2025}, can be used to prove analogous results over function fields.

We are now ready to state the main results of this paper. Our focus will be on two main examples: (i) polynomial images, and (ii) integers in a short interval. However, our method is rather general, so it should be able to flexibly handle several examples; see Section \ref{section:generalSetup} for the general but rather technical results. 

\subsection{Polynomial images} This case is particularly interesting due to its connection with conjectures of Chowla and Elliott \cite{Chowla1965, Elliott1992, Elliott1994} on correlations of multiplicative functions. See the introductions of \cite{KlurmanShkredovXu2023, ChinisShala2025} for a more detailed discussion in the setting of random multiplicative functions, where central limit theorems were proved.

\begin{theoremL}\label{thm:asboundsPoly}
    Let $f$ be a  Rademacher or Steinhaus random multiplicative function. In the Rademacher case, assume $P\in\mathbb Z[x]$ is a product of at least two distinct linear factors over $\mathbb Q$, or irreducible of degree $2$, taking infinitely many square-free values. In the Steinhaus case, assume $P(x)$ is not of the form $w(x+c)^d$ for $w\in\mathbb Z, c\in\mathbb Q$. Then almost surely we have $$\left\lvert\sum_{n\leq N} f(P(n))\right\rvert \ll \sqrt{N\log\log N}.$$ Moreover, almost surely there exists a sequence $N_k\to\infty$ such that 
    \begin{equation*}
        \left\lvert\sum_{n\leq N_k} f(P(n))\right\rvert\gg \sqrt{N_k\log\log N_k}.
    \end{equation*}
\end{theoremL}

The almost sure upper bound in Theorem \ref{thm:asboundsPoly} is new, whereas the lower bound is an extension of results of Klurman, Shkredov, and Xu \cite{KlurmanShkredovXu2023} in the Steinhaus case, and Chinis and the author \cite{ChinisShala2025} in the Rademacher case, to polynomials that split as a product of linear factors over $\mathbb Q$. Indeed, the method used in \cite{KlurmanShkredovXu2023, ChinisShala2025} relied on the fact that for a polynomial $P$ that has an irreducible factor of degree at least two, there are many integers $n$ for which $P(n)$ has a prime factor greater than $n\log n$. 

To demonstrate how this was used, let us consider $P(n) = n^2+1$. By the aforementioned fact, one can ``pull out'' a unique large prime factor for such $n$ and still capture a good ``bulk'' of the sum for infinitely many (sparse) $N$, that is 
\begin{equation*}
    \sum_{n\leq N}f(n^2+1)\approx \sum_{N\log N < p\leq N^2}f(p)\sum_{\substack{n\leq N \\ p\mid n^2+1}}f\left(\frac{n^2+1}{p}\right).
\end{equation*}

Upon conditioning on the small primes, this is a weighted sum of independent random variables, which is then handled using intricate combinatorial arguments and Gaussian approximation. 

Indeed, the method described above cannot be directly applied to the case when $P$ is a product of linear factors, for all its prime factors are of size at most linear in $N$.
In this paper, we will use a different method that does not involve any combinatorial splitting of the sum or conditioning. We will rather fully utilize the martingale structure of the entire sum -- we will elaborate on this in the next section. In particular, provided enough number-theoretic input (see the work of Martin \cite{Martin2002}), our method would also work to prove almost sure large fluctuations of $N$-smooth\footnote{An integer is called $y$-smooth if all of the prime factors dividing it are at most $y$.} values of irreducible polynomials.

\subsection{Integers in a short interval} A central limit theorem for sums of a random multiplicative function over integers in a short interval $(N-H, N]$ was first proved by Chatterjeee and Soundararajan \cite{ChatterjeeSoundararajan2012} when $H=o\left(\frac{N}{\log N}\right)$, using essentially only second and fourth moment information. This was later extended to $H\leq \frac{N}{(\log N)^{c}}$ with $c>2\log2-1$ by Soundararajan and Xu \cite{soundararajan2023clt}, entering a regime of $H$ where naively the fourth moment of the sum blows up. However, this is bypassed by throwing out $0\%$ of the integers in the short interval with atypically many prime factors, resulting in a controlled fourth moment. 

We prove an almost sure lower bound in the same regime. 

\begin{theoremL}\label{thm:asBoundsShort}
    Let $f$ be a Rademacher or Steinhaus random multiplicative function and suppose that $H=H(N)$ is a smooth, increasing, and concave function, such that $N^{\frac{11}{15}}\leq H(N)\leq N/(\log N)^c$
    for some constant $c>2\log 2-1$. 
    Then almost surely there exists a sequence $N_k\to\infty$ with corresponding $H_k = H(N_k)$ such that 
    \begin{equation*}
        \left\lvert\sum_{N_k-H_k < n \leq N_k}f(n)\right\rvert\gg \sqrt{H_k\log\frac{N_k}{H_k}}.
    \end{equation*}
\end{theoremL}

\begin{remark}
    The fluctuations of order $\sqrt{\log\frac{N}{H}}$ may seem surprisingly large at first, especially in a naive comparison with the law of iterated logarithm. However, this is explained by the fact that short sums of length $H$ decorrelate much faster, and indeed roughly $N/H$ essentially independent sums fit in a ``window'' of size $N$. In agreement with the law of iterated logarithm, the maximum of these sums has size roughly $\sqrt{\log\frac{N}{H}}$. 
\end{remark}

\begin{remark}
    The conditions on $H$ are imposed simply for convenience -- any standard choices such as $H=N^{\alpha}$ for $\alpha < 1$, $H = N/\exp\left(\sqrt{\log N}\right)$, or $H = N/(\log N)^A$ for some constant $A>2\log 2-1$, satisfy these conditions (for sufficiently large $N$). 
    
    We have not attempted to optimize the exponent $\frac{11}{15}$ -- it is likely possible to bring it down to $\frac{3}{5}$. However, the lower bound on $H$ should not be essential: for very small $H\ll \log N$ the large fluctuations are captured by almost sure long runs of ones, of length $\approx \log N$ (see the work of Erdős and Rényi \cite{ErdosRenyi1970}), whereas if $\log N\ll H\ll N^{\frac{11}{15}}$, our method should work (with significant simplifications when $H\ll N^{\frac{1}{2}}$), but the restriction stems from a result on square-free smooth integers in short intervals; see Lemma \ref{lemma:unSmoothSquarefree}. Since we are only proving a lower bound and therefore can pick convenient scales to work with, it is likely possible to circumvent this using the Matom\"aki--Radziwiłł machinery; see the work of Jain \cite{jain2025smooth} and the references within. However, we have opted for simplicity, especially since our result captures the most interesting regime where there is a transition of the large fluctuations from $\sqrt{\log N}$ down to $\sqrt{\log \log N}$.\footnote{We thank Sarvagya Jain for discussions pertaining to this remark.}
\end{remark}

It would be interesting to prove the corresponding almost sure upper bound in any regime of $H$. For rather large $H$, it is likely that the sophisticated ideas in \cite{LauTenenbaumWu2013MeanValuesRMF, Mastrostefano2022, Hardy2023AlmostSure, Atherfold2025} relating the short sum $\sum_{N-H<n\leq N} f(n)$ to a random Euler product could be useful. It is unclear, however, whether this would result in a sharp (or almost sharp) upper bound in any regime of $H$ -- this is an ongoing investigation. 

Upcoming work of Harper, Soundararajan, and Xu 
establishes a central limit theorem in the range $H = o(N)$, although with a normalization of the sum that differs from just the standard deviation. It would be interesting to determine whether our method would succeed to prove almost sure lower bounds of the right order of magnitude in this extended range. See also the work of Caich \cite{Caich2024randomshort} for more on the intricate behavior of short sums in this regime. 




\section*{Acknowledgements}
I would like to express my deep gratitude to my supervisors, Jonathan Bober and Oleksiy Klurman, for their continued support, encouragement, delightful discussions, and their careful reading and feedback on early drafts of this work. Moreover, I would like to thank Joseph Najnudel for a fruitful discussion regarding Gaussian approximation. Last but not least, I am thankful to Christopher Atherfold, Seth Hardy, and Neo Tardy, for enlightening technical conversations and their comments on earlier drafts of the paper.  The author is funded by a University of Bristol PhD scholarship.

\section{Overview of the proofs}
\subsection{Martingale structure} Our starting point is the observation of Harper \cite{Harper2013} that sums of a random multiplicative function have the structure of a martingale with respect to the filtration induced by the largest prime factor, namely 
\begin{equation*}
    \sum_{n\in\mathcal A}f(n) = \sum_{p^{\alpha}\leq N}f(p^{\alpha})\sum_{\substack{n\in\mathcal A \\ P^+(n)=p \\ p^{\alpha}\parallel n}}f\left(\frac{n}{p^{\alpha}}\right).
\end{equation*}
This observation was utilized in almost all subsequent works on random multiplicative functions. For us, its usefulness lies in the fact that it morally reduces the understanding of $\sum_{n\in\mathcal A}f(n)$ to evaluating the second and fourth moments of the sum, provided they do not blow up. More precisely, we will apply the following quantitative result from the book of Hall and Heyde \cite{HallHeyde1980}, yielding a bound on the Kolmogorov distance between our martingale and a Gaussian. 
\begin{prop}[\cite{HallHeyde1980}, Theorem 3.9]\label{prop:quantCLT}
    Let $X_1, X_2, \ldots, X_n$ be a real square-integrable martingale difference sequence.
    %with respect to the filtration $\mathcal F_1\subseteq\mathcal F_2 \subseteq\cdots\subseteq\mathcal F_n$. 
    Let $S_n = \sum_{i=1}^n X_i$. Then 
    \begin{equation*}
\left\lvert\mathbb P(S_n\leq x) - \frac{1}{\sqrt{2\pi}}\int_{-\infty}^x e^{-t^2/2}\emph{d}t\right\rvert\ll \frac{\left(\sum_{i=1}^n \mathbb E\lvert X_i\rvert^4 + \mathbb E\left\lvert\sum_{i=1}^n X_i^2 - 1\right\rvert^2\right)^{1/5}}{1+\lvert x\rvert^{\frac{16}{5}}}.
    \end{equation*}
\end{prop}
\subsection{Almost sure upper bound} Equipped with the above result, we may readily prove an almost sure upper bound for $\sum_{n\in\mathcal A} f(n)$, provided our approximation is suitably strong. In order to access the tails at $x\approx C\sqrt{\log k}$ which occur with probability $\approx k^{-C^2/2}$ for the Gaussian, we will need an approximation that is at least this good. We are able to achieve this in the case of polynomial images, but not short intervals -- we will discuss this in more detail in the following subsections. Using our quantitative approximation at several spread out scales together with a slow variation property between the scales (see Proposition \ref{prop:slowVar} and Lemma \ref{lemma:slowVarPoly}), the Borel-Cantelli lemma establishes the almost sure upper bound. 

\subsection{Almost sure lower bound} Establishing an almost sure lower bound is where the main novelty of this paper lies. The goal here is obtaining a quantitative multivariate central limit theorem for 
\begin{equation*}
    \left(\frac{1}{\sqrt{\lvert\mathcal A_1\rvert}}\sum_{n\in\mathcal A_1}f(n), \frac{1}{\sqrt{\lvert\mathcal A_2\rvert}}\sum_{n\in\mathcal A_2}f(n), \ldots, \frac{1}{\sqrt{\lvert\mathcal A_k\rvert}}\sum_{n\in\mathcal A_k}f(n) \right), 
\end{equation*}
where each $\mathcal A_l\subseteq [1, N_l]\cap \mathbb N$ for various integers $N_1, N_2, \ldots, N_k$. To do this, roughly speaking, we will first apply Proposition \ref{prop:quantCLT} to linear combinations of the normalized sums $\lvert\mathcal A_l\rvert^{-\frac{1}{2}}\sum_{n\in\mathcal A_l} f(n)$, establishing approximate joint Gaussianity. We will then bootstrap this to a quantitative multivariate central limit theorem, by using the following recent key result of Bobkov and G\"otze \cite{BobkovGoetze2024} to bound the Kantorovich--Wasserstein distance between the multivariate distributions in terms of the Kolmogorov distances of the corresponding linear combinations. 

\begin{prop}[\cite{BobkovGoetze2024}, Theorem 1.1]\label{prop:multivarCLT} Let $X$ and $Y$ be $\mathbb R^k$-valued random variables with $\|X\|_p\leq b$ and $\|Y\|_p\leq b$ for some $p>1$ and $b\geq 0$, where $\|Z\|_p = \left( \mathbb E\lvert Z\rvert^p\right)^{\frac{1}{p}}$. Then for $q$ such that $1/p + 1/q=1$, we have
\begin{multline*}
    \sup_{\|u\|_{\emph{Lip}}\leq 1} \lvert \mathbb EuX - \mathbb EuY\rvert \\ \ll b^{1-2/(qk+2)} \left(\sup_{\lvert \theta\rvert = 1} \int_{-\infty}^\infty \left\lvert\mathbb P\left(\sum_{l=1}^k \theta_l X_l\leq x \right)- \mathbb P\left(\sum_{l=1}^k \theta_l Y_l\leq x\right)\right\rvert\text{d} x\right)^{2/(kq+2)}. 
\end{multline*} Here the supremum is taken over all functions $u:\mathbb R^k\to\mathbb R$ with Lipschitz semi-norm at most $1$. 
\end{prop}

This is a more direct approach compared to the use of Stein's method via exchangeable pairs in the work of Harper \cite{Harper2013ZetaMax, Harper2023_ASLF}. In particular, the author was unable to adapt the latter method to handle the complex dependencies between large primes $p$ and $q$ such that $p\mid n$ and $q\mid n+1$ when dealing with $\sum_{n\leq N} f(n(n+1))$, say. The advantage of our approach is that we obtain a quantitative multivariate central limit theorem purely from the one-dimensional martingale structure. See also the recent work of Kowalski and Untrau \cite{KowalskiUntrau2025Wasserstein}, Humphries \cite{Humphries2025QuantitativeEquidistribution}, as well as earlier work of Saksman and Webb \cite{saksman2020riemann} on the use of the Kantorovich--Wasserstein distance in number theory. 

Equipped with a quantitative multivariate central limit theorem, we may compare the distribution of our $k$ sums to the distribution of $k$ Gaussian random variables with a prescribed covariance structure governed by the arithmetic of the set $\mathcal A$. Given this information, we may readily establish the almost sure lower bound by an application of Proposition \ref{prop:multivarCLT} and the Borel-Cantelli lemma. 

\subsection{Arithmetic input} After applying Proposition \ref{prop:quantCLT} to $\sum_l\theta_l\lvert \mathcal A_{N_l}\rvert^{-\frac{1}{2}}\sum_{n\in\mathcal A_l}f(n)$ for Rade-macher (or Steinhaus) $f$, what we essentially have to control is the number of solutions to the equation $$n_1n_2n_3n_4=\square \text{ (or } n_1n_2=n_3n_4\text{)}$$ with $n_i\in\mathcal A_{N_{l_i}}$ for $1\leq i\leq 4$ coming from (at most) four different scales. Henceforth we will refer to this as the \emph{fourth moment equation}. 
%It is helpful to keep in mind that if we allow $n_i$ to range through square-free (all) integers up to $N_{l_i}$ and the scales $N_l$ are not too far apart (let us say $\log N_l\sim \log X$ for all $1\leq l\leq k$), the number of solutions to the equation is approximately $\sqrt{N_{l_1}N_{l_2}N_{l_3}N_{l_4}}(\log X)^2$ (or $\sqrt{N_{l_1}N_{l_2}N_{l_3}N_{l_4}} \log X$ in the Steinhaus case).
In each of the examples we consider in the paper, we would like for this equation to ``only'' (in an asymptotic sense) have the trivial solutions where the $n_i$ are equal in pairs, contributing approximately $N_{l_i}N_{l_j}$ solutions for $1\leq i, j\leq 4$. This means that we would like to rule out the many non-trivial solutions that occur in the full ranges of integers. In the following subsections, we will describe how we achieve this, as well as the subtleties in the previously described method for the examples we consider in this paper. 

\subsubsection{Polynomial images} This case will be our most streamlined application of the method as described above. 
We will utilize the main results of \cite{KlurmanShkredovXu2023, ChinisShala2025}, giving a power-saving bound on the number of non-trivial solutions to the fourth moment equations  $$P(n_1)P(n_2)P(n_3)P(n_4) = \square \text{ (or  } P(n_1)P(n_2) = P(n_3)P(n_4)\text{)}$$ with all $n_i\leq N$. This allows us to choose the $N_l$ for $1\leq l\leq k$ so that we can control the covariances of the sums $N_l^{-\frac{1}{2}}\sum_{n\leq N_l} f(P(n))$ quite crudely, by simply bounding the number of non-trivial solutions to the equation with $n_i$ ranging up to the maximum of the $N_{l_i}$. 

However, the martingale structure that we use forces us to combine the power-saving bound above with upper bounds for (very) smooth values of polynomial images, resulting in a quantitative bound of the form $\exp\left(-\sqrt{\log X}\right)$ in the martingale central limit theorem (where $X$ is so that $\log N_l\asymp \log X$ for all $1\leq l\leq k$); see Corollary \ref{cor:quantCLTpoly}. Nonetheless, this is certainly strong enough to allow us to access the tails that occur with probability $\approx (\log X)^{-C^2/2}$, leading to the almost sure upper bound as described before.

\begin{remark}\label{rmk: wangXumoments} One could also obtain the almost sure upper bound in the Steinhaus case by evaluating the $2k$-th  moment of $\sum_{n\leq N} f(P(n))$ for $k\approx \log\log N$, as was done in the work of Wang and Xu \cite{Wang-Xu}. A careful inspection of their proof shows that this is indeed the uniformity they obtain. This appears harder to do in the Rademacher case.\footnote{We thank Christopher Atherfold for discussions pertaining to this remark.}  \end{remark}

For the lower bound, passing to the multivariate central limit theorem will result in a further loss, namely a bound of the form $\sqrt k\exp\left(-\frac{\sqrt{\log X}}{k}\right)$ for the approximation; see Corollary \ref{cor:MCLTpoly}. However, it will be sufficient to choose $k=(\log X)^{\varepsilon_0}$ for some small $\varepsilon_0>0$ in order to create large fluctuations of order $\approx \sqrt{2\log k}\approx \sqrt{\log\log X}$ almost surely. 

\subsubsection{Integers in a short interval} Attempting to follow the strategy as in the polynomial case here runs into the following difficulties. Firstly, the bounds we have for the number of non-trivial solutions to $n_1n_2n_3n_4=\square$ (or $n_1n_2=n_3n_4$ in the Steinhaus case) are much weaker and get worse as $H$ gets bigger. Secondly, even when $H=N^{\alpha}$ in which case we have power-saving bounds on the number of non-trivial solutions, as before, the very smooth solutions (including the trivial ones) yield a martingale central limit theorem with a quantitative bound of shape roughly $\exp\left(-\sqrt{\log X}\right)$. This is not enough to capture the large fluctuations of size $C\sqrt{\log X}$ that occur probability roughly $X^{-C^2/2}.$ 

\begin{remark}
    When $k=1$, the above strategy nonetheless yields a significant quantitative improvement of the central limit theorem of Chatterjee and Soundararajan \cite{ChatterjeeSoundararajan2012}, but we do not state this in the paper. We note that the result of Soundararajan and Xu \cite{soundararajan2023clt} is also quantitative, but stated in a different form in terms of the characteristic function, with an apparent large loss of $e^{t^2/2}$. However, it appears that by being less crude with the bound $|1+it|\leq e^{t^2/2}$ for large $t$ in their work, applying Esseen's inequality (truncated Fourier inversion) to their result would probably recover a quantitative central limit theorem of similar strength. 
\end{remark}

To bypass these issues, we condition on the values of $f$ on the small primes, and approximate the sums $\sum_{N_l-H_l<n\leq N_l} f(n)$ for most $1\leq l\leq k$ with a weighted sum of independent random variables $f(p)$ for large primes $p$. This will be advantageous, as the weaker saving we have on the number of non-trivial solutions to the fourth moment equation will only be used to control the $\approx k^2$ conditional covariances, and will not play a role in the quantitative strength of the conditional martingale central limit theorem. 


In controlling the conditional covariances by counting solutions to the equation $n_1n_2n_3n_4 = \square$ (or $n_1n_2=n_3n_4$) where the $n_i$ are potentially coming from different scales, we have to be mindful about the choice of scales. This is due to the many non-trivial solutions (of the same order of magnitude as the trivial solutions) which are obtained by scaling, such as solutions of the form $(n_1)(2n_1)(n_3)(2n_3)=\square$, say. 

When $H$ is greater than roughly $\frac{N}{\log N}$, we implement the strategy used in the work of Soundararajan and Xu \cite{soundararajan2023clt}, throwing out integers with atypically many prime factors, resulting in the desired control of the fourth moment equation as long as $H\leq \frac{N}{(\log N)^c}$ with $c>2\log 2-1$; see Proposition \ref{prop:CountAt4Scales}. For us, however, it is crucial that we obtain a sufficiently strong quantitative bound for the number of such integers; see Lemmas \ref{lemma:tooManyPrimes} and  \ref{lemma:ignoreTooManyPrimesShort}.



\subsubsection{Other examples and limitations} The method used in this paper may be used to prove almost sure lower bounds on large fluctuations of other examples, such as: shifted primes, sums of squares in short intervals \cite{soundararajan2023clt}, rough integers \cite{Xu2024btran}, and exponential sums with random multiplicative coefficients $\sum_{n\leq N} f(n)e(n\theta)$ \cite{BenatarNishryRodgers2022, Hardy2024imrn, soundararajan2023clt}. However, we are not able to handle the case of restricted number of prime factors $\omega(n)\leq w=o(\log\log N)$ for $n\le N$ considered in \cite{Hough2011, Harper2013}, for the number of non-trivial solutions to the fourth moment equation there is very large. Even when $w$ is finite, the saving over the number of trivial solutions is at most double-logarithmic in $N$, resulting in a very small choice for the number of scales $k$ that we can hope to use. It would be interesting to determine whether this is a reflection of the truth, that is, whether the large fluctuations of $\sum_{n\leq N,\text{ } \omega(n)\leq w} f(n)$ differ from what would be expected from the approximate Gaussian distribution. 

In another direction, it would also be interesting to prove almost sure upper bounds for other examples, such as shifted primes. 
In contrast with short intervals, say, it is not clear how one would proceed here, as this sum likely cannot be related to a random Euler product --- this is an ongoing investigation. 

\section{Auxiliary Results}

In this section, we collect and/or prove some auxiliary results that we will use throughout our proofs.

\subsection{Number theory results} We start by stating a result of Nair and Tenenbaum \cite{NairTenenbaum1998} (see also the work of Henriot \cite{Henriot2012_NairTenenbaumUniform}), in a less general form that will be sufficient for us. This is a generalization of a well-known lemma of Shiu. 

\begin{prop}\label{prop:nair-tenenbaum}
    Let $F$ be a non-negative multiplicative function such that there exists a constant $A\geq 1$ and $\epsilon >0$ such that $F(m)\leq Am^{\epsilon}$ for all $m\in\mathbb N$. Let $Q\in\mathbb Z[x]$ be a fixed polynomial of degree $d$ with no fixed prime factor, and let $\rho(n)$ be the number of solutions to $Q(x)\equiv 0\pmod n$. Suppose that $Q$ has $r$ distinct irreducible factors $Q_1, Q_2, \ldots, Q_r$, and write $Q = \prod_{i=1}^r Q_i^{\gamma_i}$. Let $\rho_i(n)$ be the number of solutions to $Q_i(x)\equiv 0\pmod n$ for $1\leq i\leq r$. Denote the discriminant of the square-free kernel of $Q$ by $D$. If $\epsilon<\frac{1}{8d^2}$, then 
    \begin{multline*}
        \sum_{N<n\leq N+H} F(\lvert Q(n)\rvert ) \\ \ll H\prod_{p\leq N}\left(1-\frac{\rho(p)}{p}\right) \sum_{\substack{n_1^{\gamma_1}n_2^{\gamma_2}\cdots n_r^{\gamma_r}\leq N \\ (n_i, n_j)=1 \text{ }\forall i\neq j \\ (n_i, D)=1 \text{ }\forall i}} F(n_1^{\gamma_1}\cdots n_r^{\gamma_r})\frac{\rho_1(n_1)\rho_2(n_2)\cdots \rho_r(n_r)}{n_1n_2\cdots n_r},
    \end{multline*}
    whenever $N$ is large enough, and $N^{4d^2\epsilon}\leq H\leq N$. 
\end{prop}

We now record a few corollaries of this result. 

\begin{lemma}\label{lemma:PolyBound}
    Let $Q\in\mathbb Z[x]$ be a polynomial. There exists a constant\footnote{In fact Proposition \ref{prop:nair-tenenbaum} gives a much more precise description of $c_Q$, but we will not need it here.} $c_Q>0$ such that 
    \begin{equation*}
        \sum_{N<n\leq N+H} \tau_3(\lvert Q(n)\rvert)\ll H (\log N)^{c_Q}
    \end{equation*}
    whenever $N^{1/3}\leq H \leq N$, where $\tau_3$ is the $3$-fold divisor function.
\end{lemma}

\begin{proof}
We have $\tau_3(n)\ll_{\epsilon} n^{\epsilon}$ for every $\epsilon>0$. Suppose $Q$ has $r$ distinct irreducible factors with the largest degree of an irreducible factor being $d$. Applying Proposition \ref{prop:nair-tenenbaum} to $\tau_3$ and the polynomial $Q$, and using the submultiplicativity of $\tau_3$, we obtain 
\begin{align*}
     \sum_{N<n\leq N+H} \tau_3(\lvert Q(n)\rvert) & \ll H (\log N)^{-\deg Q}\sum_{n_1\cdots n_r\leq N}\tau_3(n_1\cdots n_r)\frac{\rho_1(n_1)\cdots \rho_r(n_r)}{n_1\cdots n_r}\\ &  \ll H(\log N)^{-\deg Q} \left(\exp\left(\sum_{p\leq N} \frac{3d}{p} \right)\right)^r \ll H(\log N)^{3dr - \deg Q},
\end{align*} whenever $N\geq H\geq N^{4\deg^2 Q\epsilon}\geq N^{\frac{1}{3}}$ if $\epsilon$ is chosen small enough. 
\end{proof}

\begin{lemma}\label{lemma:tooManyPrimes}
    Let $\varepsilon>0$ and let $N$ and $H$ be large enough such that $H\geq N^{5\varepsilon}$. Let $\Omega(n)$ denote the number of prime factors of $n$, counted with multiplicity. Then there exists $\varepsilon'>0$ such that 
    \begin{equation*}
        \#\{N-H<n\leq N : \Omega(n)>(1+\varepsilon)\log\log N\}\ll \frac{H}{(\log N)^{\varepsilon'}}.
    \end{equation*}
\end{lemma}

\begin{proof}
The quantity we want to bound is certainly at most 
\begin{equation*}
    \frac{1}{\exp((1+\varepsilon)\log\log N\log(1+\varepsilon))}\sum_{N-H<n\leq N} (1+\varepsilon)^{\Omega(n)}.
\end{equation*}
Applying Proposition \ref{prop:nair-tenenbaum} with  $A=1, \epsilon = \log(1+\varepsilon)/\log 2=\varepsilon/\log 2 +O(\varepsilon^2)$ to the sum gives 
\begin{equation*}
    \sum_{N-H<n\leq N} (1+\varepsilon)^{\Omega(n)}\ll H(\log N)^{\varepsilon} = H\exp(\varepsilon\log\log N),
\end{equation*}
whenever $H\geq N^{4\epsilon}\gg N^{5\varepsilon}$. Putting the two bounds together gives that the quantity we want to bound is $\ll H/(\log N)^{(1+\varepsilon)\log(1+\varepsilon)-\varepsilon}$, so we may take $\varepsilon' = (1+\varepsilon)\log(1+\varepsilon)-\varepsilon = \frac{\varepsilon^2}{2} + O(\varepsilon^3)$. 
\end{proof}

\begin{lemma}\label{lemma:unSmoothSquarefree}
    Let $N$ be large, $\alpha>\frac{1}{2}$, and let $N\geq H=H(N)\geq N^{\frac{4\alpha + 1}{5}}$. The number of square-free integers in the short interval $(N-H, N]$ such that $P^+(n)>N^{\alpha}$ is 
    \begin{equation*}
        \gg H+O\left(\frac{H}{\log N} \right).
    \end{equation*}
\end{lemma}

\begin{proof}
    The number of such $n$ is at least
    \begin{equation*}
        \sum_{N^{\alpha} < p\leq (H^5/N)^{1/4}}\sum_{\frac{N-H}{p}<n\leq \frac{N}{p}}\mu^2(n).
    \end{equation*}
    According to (a weaker version of) \cite[Theorem 1.1]{Pandey2024Squarefree}, for the inner sum we obtain 
    \begin{equation*}
        \sum_{\frac{N-H}{p}<n\leq \frac{N}{p}}\mu^2 (n) = \frac{6}{\pi^2}\frac{H}{p} + O\left(\frac{N^{\frac{1}{5}}}{p^{\frac{1}{5}}} \right).
    \end{equation*}
    Summing over the primes $N^{\alpha}<p\leq (H^5/N)^{1/4}$, we obtain 
    \begin{equation*}
        \gg H+O\left(H/\log N \right),
    \end{equation*}
giving the desired conclusion. 
\end{proof}



\subsection{Probability results} We begin by adapting Proposition \ref{prop:quantCLT} to the complex case for the Steinhaus model. 

\begin{prop}\label{prop:quantCLT-complex}
    Let $X_1, X_2, \ldots, X_n$ be a complex square-integrable martingale difference sequence. 
    %with respect to the filtration $\mathcal F_1\subseteq\mathcal F_2 \subseteq\cdots\subseteq\mathcal F_n$. 
    Let $S_n = \sum_{i=1}^n X_i$. Then 
    \begin{multline*}
\left\lvert\mathbb P(\mathfrak TS_n\leq x) - \frac{1}{\sqrt{2\pi}}\int_{-\infty}^xe^{-\frac{t^2}{2}}\emph{d}t\right\rvert \\ \ll \frac{\left(\sum_{i=1}^n \mathbb E\lvert X_i\rvert^4 + \mathbb E\left\lvert\sum_{i=1}^n \lvert X_i\rvert ^2 - 1\right\rvert^2 + \mathbb E\left\lvert\sum_{i=1}^n \left(X_i^2 +\overline{X_i^2}\right)\right\rvert^2\right)^{1/5}}{1+\lvert x\rvert^{\frac{16}{5}}},
    \end{multline*} 
    where $\mathfrak T$ denotes $\sqrt2$ times either the real or imaginary part of $S_n$. 
\end{prop}

\begin{proof}
    This follows easily from applying Proposition \ref{prop:quantCLT} to the martingales $(S_n \pm \overline{S_n})/2$ and using the inequalities
    \begin{equation*}
        \left\lvert X_i \pm \overline{X_i}\right\rvert^4 \ll \lvert X_i\rvert^4 \hspace{0.5cm}\text{ and }\hspace{0.5cm} \left\lvert \left(X_i \pm \overline{X_i}\right)^2 - 1 \right\rvert^2 \ll \left\lvert \lvert X_i\rvert^2 - 1\right\rvert^2 + \left\lvert X_i^2  + \overline{X_i^2} \right\rvert^2.
    \end{equation*}
More details in a related setting are provided in the proof of Theorem 2.4 in \cite{soundararajan2023clt}. 
\end{proof}

Next, we borrow a result from Billingsley's book \cite{Billingsley1999Convergence}, bootstrapping fourth moment information in short intervals to a deviation estimate for the maximum of a sum of random variables. 

\begin{lemma}[Theorem 10.2 in \cite{Billingsley1999Convergence}] \label{lemma:billings} Let $S_N = X_1 + X_2 + \ldots + X_N$ be a sum of random variables. Let $u_1, u_2, \ldots, u_n$ be real numbers such that 
\begin{equation*}
    \mathbb P(\lvert S_j - S_i\rvert\geq\lambda)\leq\frac{1}{\lambda^4}\left(\sum_{i<k\leq j} u_k\right)^2 
\end{equation*} for all $1\leq i<j\leq n$. Then 
\begin{equation*}
    \mathbb P\left(\max_{k\leq n} \lvert S_k\rvert\geq\lambda \right) \ll \frac{1}{\lambda^4}\left(\sum_{1\leq k\leq n} u_k\right)^2 .
\end{equation*}
    
\end{lemma}

To prove the almost sure lower bounds, we will actually show that our sums attain large values with large probability rather frequently. The next lemma establishes that this is sufficient.  

\begin{lemma}\label{lemma:asLbGeneral}
Let $\mathcal A_N$ be a collection of subsets of positive integers, and let $X_n$ for $n\in\mathcal A_N$ be random variables. Let $\mathcal L(X)$ and $\mathcal E(X)\to_{X\to \infty} 0$ be functions such that $\mathcal L(x)\asymp \mathcal L(X)$ for all $X<x\leq X^2$, and 
\begin{equation*}
    \max_{X<N\leq X^2}\left\lvert\sum_{n\in\mathcal A_N} X_n \right\rvert\geq \mathcal L(X)
\end{equation*} for every large enough $X$ with probability $1-\mathcal E(X)$. Then there exists a sequence $N_l\to\infty$ such that almost surely
\begin{equation}\label{eq:LargeL}
    \left\lvert \sum_{n\in\mathcal A_{N_l}} X_n \right\rvert \gg \mathcal L(N_l).
\end{equation}
\end{lemma}

\begin{proof}
   Begin with a sparse sequence $N_l$ satisfying $N_{l+1}>N_l^2$, say. For a fixed $l$, the probability that \eqref{eq:LargeL} fails for all $N\in (X_l, X_l^2]$ and a suitably small implied constant is at most $\mathcal E(X_l)$ by the condition of the proposition. Now we may choose a subsequence $N_{l_k}$ of $N_l$ so sparse depending on $\mathcal E$ such that $\sum_{k=1}^\infty \mathcal E\left(N_{l_k}\right)$ converges. Then by Borel-Cantelli, \eqref{eq:LargeL} fails for at most finitely many $l_k$. 
\end{proof}

Finally, we record a quantitative version of Slepian's lemma. 

\begin{lemma}[\cite{Harper2023_ASLF}, Normal Comparison Result 1]\label{lemma:normalsBigLemma} Suppose that $k \geq 2$, and that $\epsilon \geq 0$ is sufficiently small (i.e., less than a certain small absolute constant). Let $Y_1, Y_2, \dots, Y_k$ be
mean zero, variance one, jointly normal random variables, and suppose $\mathbb EY_{l_1}Y_{l_2}\leq \epsilon$ whenever $l_1\neq l_2$. Then, for any $100\epsilon \leq \delta \leq 1/100$, we have
\[
\mathbb{P}\left(\max_{1\leq l\leq k} Y_l\leq \sqrt{(2-\delta)\log k}\right) \leq  \exp\left(-\Theta\left(\frac{k^{\delta/20}}{\sqrt{\log k}}\right)\right) + k^{-\delta^2/(50\epsilon)}.
\]
\end{lemma}


\section{General Framework}\label{section:generalSetup}
We are going to deduce our results from the following general framework. Throughout the rest of this paper, let $f$ be a Rademacher (or Steinhaus) random multiplicative function. We are interested in subsets $\mathcal A_N\subseteq\{1, 2, \ldots, N\}$ of square-free (or all) positive integers, sampled at various $N_1, N_2, \ldots, N_k$, with (some of) the following properties for any $1\leq l_1, \cdots , l_4\leq k$:
\begin{itemize}
    \item[(I)] There exists $\varepsilon_1>0$ such that the number of non-trivial\footnote{We call a solution \emph{non-trivial} if the $n_i$ are not equal in pairs.} solutions to $$n_1n_2n_3n_4=\square \text{ }(\text{or }n_1n_4=n_2n_3)$$ 
    with $n_i\in\mathcal A_{N_{l_i}}$ and $P^+(n_1)=P^+(n_2), P^+(n_3)=P^+(n_4)$ is $\leq \varepsilon_1\sqrt{\lvert\mathcal A_{N_{l_1}}\rvert\cdots \lvert\mathcal A_{N_{l_4}}\rvert}$.
     \item[(I')] There exists $\varepsilon_1'>0$ such that the number of non-trivial solutions to $n_1n_2n_3n_4=\square$ 
     (or $n_1n_4=n_2n_3$) 
    with $n_i\in\mathcal A_{N_{l_i}}$ and $P^+(n_1) = \cdots = P^+(n_4)$ is $\leq \varepsilon_1'\sqrt{\lvert\mathcal A_{N_{l_1}}\rvert\cdots \lvert\mathcal A_{N_{l_4}}\rvert}$. 
    
    (Note that $\varepsilon'_1$ may be taken to be at most $\varepsilon_1$ if (I) is satisfied.)
    
  
    
    \item[(II)] There exists $\varepsilon_2>0$ such that the number of $n_1\in\mathcal A_{N_{l_1}}$ and $n_2\in\mathcal A_{N_{l_2}}$ such that $P^+(n)=P^+m$ is $\leq \varepsilon_2\lvert\mathcal A_{N_{l_1}}\rvert\lvert\mathcal A_{N_{l_2}}\rvert$.


    \item[(III)] We either have: 
    \begin{itemize}
        \item[(III$_a$)] If $N_{l_1}\leq N_{l_2}$ then $\mathcal A_{N_{l_1}}\subseteq \mathcal A_{N_{l_2}}$, or
        \item[(III$_b$)] If $l_1\neq l_2$, then $\mathcal A_{N_{l_1}}\cap\mathcal A_{N_{l_2}} =\emptyset$. 
    \end{itemize} 

\end{itemize}



\begin{remark}
    Condition (III) is automatic if we begin with an infinite subset of square-free (or all) positive integers $\mathcal A$, and let $\mathcal A_N = \mathcal A \cap(N-H, N]$ for some $H=H(N)\to\infty$, as long as the scales $N_1, N_2, \ldots, N_k$ are sufficiently separated. 
\end{remark}
Let
\begin{equation}\label{eq:S_Ndef}
    S_N = \frac{1}{\sqrt{\lvert \mathcal A_N\rvert}}\sum_{n\in\mathcal A_N} \mathfrak Tf(n) = \sum_{p} \frac{1}{\sqrt{\lvert \mathcal A_N\rvert}} \sum_{\substack{n\in\mathcal A_N \\ P^+(n) = p}} \mathfrak Tf(n) =: \sum_{p} M_{p, N},
\end{equation}
where here and throughout the paper $\mathfrak T$ denotes simply the identity in the Rademacher case, and $\sqrt2$ times either the real or imaginary part in the Steinhaus case. 

\subsection{Approximate joint Gaussianity}
For any sequence of integers $N_1< N_2< \cdots< N_k$ and real coefficients $c_1, c_2, \ldots, c_k$ with $\sum_{l=1}^k \lvert c_l\rvert^2 = 1$, the sum 
\begin{equation*}
    c_1S_{N_1} + c_2S_{N_2} + \cdots + c_kS_{N_k} = \sum_p \sum_{l=1}^k c_lM_{p, N_l}
\end{equation*} is a martingale, with $\sum_{l=1}^k c_lM_{p, N_l}$ being a martingale difference sequence indexed by primes $p$. 


Applying Proposition \ref{prop:quantCLT} (or Proposition \ref{prop:quantCLT-complex}) to $\sum_{l=1}^k c_lM_{p, N_l}$ leads us to bounding the quantities 
\begin{equation*}
    A := \sum_{p} \mathbb E\left\lvert\sum_{l=1}^k c_lM_{p, N_l}\right\rvert^4 \hspace{1cm}\text{and}\hspace{1cm} B:= \mathbb E\left\lvert\sum_{p} \left\lvert \sum_{l=1}^k c_lM_{p, N_l}\right\rvert^2 - 1\right\rvert^2.
\end{equation*}
In the Steinhaus case, for our particular martingale, the last term from Proposition \ref{prop:quantCLT-complex} is $\ll A$. This is because upon expanding
\begin{equation*}
    \mathbb E \left\lvert\sum_p \left(\left(\sum_{l=1}^k c_lM_{p, N_l}\right)^2 + \left(\sum_{l=1}^k \overline{c_lM_{p, N_l}}\right)^2\right) \right\rvert^2,
\end{equation*}
for all primes $p\neq q$ we have $$\mathbb EM_{p, N_{l_1}}^2M_{q, N_{l_2}}^2=\mathbb E\overline{M_{p, N_{l_1}}^2}M_{q, N_{l_2}}^2 = \mathbb EM_{p, N_{l_1}}^2\overline{M_{q, N_{l_2}}^2} = \mathbb E\overline{M_{p, N_{l_1}}^2M_{q, N_{l_2}}^2}=0,$$ whereas for the primes $p=q$ the surviving terms are (note that $\mathbb E\left(\sum_{l=1}^k c_l M_{p, N_l}\right)^4 =$ \linebreak $\mathbb E\overline{\left(\sum_{l=1}^k c_l M_{p, N_l}\right)}^4 =  0$ since $\mathbb Ef(p)^4 = 0$)
\begin{equation*}
    \ll \sum_{p} \mathbb E\left(\sum_{l=1}^k c_lM_{p, N_l}\right)^2\overline{\left(\sum_{l=1}^k c_l M_{p, N_l}\right)^2} = A.
\end{equation*}

\subsubsection{Bounding $A$} We have
\begin{equation*}
    A = \sum_p\sum_{1\leq l_i\leq k} c_{l_1}\cdots \overline{c}_{l_4}\mathbb E M_{p, N_{l_1}}\cdots \overline{M}_{p, N_{l_4}}. 
\end{equation*}
Further, 
\begin{equation*}
    \mathbb E M_{p, N_{l_1}}\cdots \overline{M}_{p, N_{l_4}} = \frac{1}{\sqrt{\lvert \mathcal A_{N_{l_1}}\rvert\cdots \lvert \mathcal A_{N_{l_4}}\rvert}}\sum_{\substack{n_i\in\mathcal A_{N_{l_i}} \\ P^+(n_i) = p}} \mathbb Ef(n_1)\cdots \overline{f(n_4)}.
\end{equation*}

Note that $\mathbb Ef(n_1)\cdots \overline{f(n_4)}$ is non-zero only if $n_1n_2n_3n_4$ is a perfect square (or $n_1n_2=n_3n_4$), in which case the expectation is equal to $1$. Therefore, by moving the sum over $p$ inside, we have 
\begin{equation*}
    A = \sum_{1\leq l_i\leq k} \frac{c_{l_1}}{\sqrt{\lvert\mathcal A_{N_{l_1}}\rvert}}\cdots \frac{\overline{c}_{l_4}}{\sqrt{\lvert\mathcal A_{N_{l_4}}}\rvert}\sum_{\substack{n_i\in\mathcal A_{N_{l_i}} \\ n_1n_2n_3n_4 = \square \text{ }(\text{or } n_1n_2=n_3n_4) \\ P^+ (n_i) = P^+ (n_j) \text{ }\forall i, j }}1.
\end{equation*}
The contribution from the trivial solutions (note that by condition (III), $n_1=n_2$ forces one of $\mathcal A_{N_{l_1}}, \mathcal A_{N_{l_2}}$ to be a subset of the other, and similarly for the other variables) is 
\begin{equation*}
\ll  \sum_{\substack{1\leq l_i\leq k }} \frac{\lvert c_{l_1}\rvert\cdots \lvert c_{l_4}\rvert}{\sqrt{\lvert \mathcal A_{N_{l_1}}\rvert\cdots \lvert\mathcal A_{N_{l_4}}\rvert}}\sum_{\substack{n_1\in\mathcal A_{N_{l_1}}\cap \mathcal A_{N_{l_2}} \\ n_3\in\mathcal A_{N_{l_3}}\cap \mathcal A_{N_{l_4}} \\ P^+(n_1) = P^+(n_3)}} 1\ll \varepsilon_2k^2, 
\end{equation*}
by condition (II). Here we used that if, say $\mathcal A_{N_{l_1}}\cap \mathcal A_{N_{l_2}} = \mathcal A_{N_{l_1}}$, then certainly $\lvert\mathcal A_{N_{l_2}}\rvert\geq\lvert\mathcal A_{N_{l_1}}\rvert$, and then applied the Cauchy-Schwarz inequality for the sum over the $l_i$. Note that if (III$_b$) is satisfied, then $n_1=n_2$ forces $l_1=l_2$ (and similarly for the other variables), in which case the contribution from the trivial solutions would be $\ll\varepsilon_2 \sum_{1\leq l_1, l_3\leq k} \lvert c_{l_1}\rvert^2\lvert c_{l_3}\rvert^2\ll \varepsilon_2.$

The contribution from non-trivial solutions is
\begin{equation*}
    \ll\varepsilon_1' \sum_{1\leq l_i\leq k} \lvert c_{l_1}\rvert\cdots \lvert c_{l_4}\rvert \ll \varepsilon_1'k^2 
\end{equation*}
using $\sum_{l=1}^k \lvert c_{l}\rvert\ll\sqrt k$ by Cauchy-Schwarz and condition (I'). 


We conclude that
\begin{equation}\label{eq:Abound}
    A \ll k^2(\varepsilon_1'  + \varepsilon_2),
\end{equation}
or 
\begin{equation}\label{eq:AboundDisjoint}
    A\ll k^2\varepsilon_1' + \varepsilon_2
\end{equation}
if (III$_b$) is satisfied. 




\subsubsection{Bounding $B$} We have 
\begin{multline*}
    B = 1+ \sum_{p, q} \sum_{1\leq l_i\leq k} \frac{c_{l_1}}{\sqrt{\lvert \mathcal A_{N_{l_1}}\rvert}}\cdots \frac{c_{l_4}}{\sqrt{\lvert \mathcal A_{N_{l_4}}\rvert}}\sum_{\substack{n_i\in\mathcal A_{N_{l_i}} \\ n_1n_2n_3n_4 = \square \text{ }(n_1n_4=n_2n_3) \\ P^+(n_1) = P^+(n_2)=p \\ P^+(n_3) = P^+(n_4) =q}} 1  \\ \\ - 2\sum_p\sum_{1\leq l_i\leq k} \frac{c_{l_1}}{\sqrt{\lvert \mathcal A_{N_{l_1}}\rvert}}\frac{\overline{c}_{l_2}}{\sqrt{\lvert \mathcal A_{N_{l_2}}\rvert}}\sum_{\substack{n_i\in\mathcal A_{N_{l_i}} \\ n_1n_2=\square\text{ }(\text{or }n_1=n_2) \\ P^+(n_1) = P^+(n_2) = p}} 1 \label{Beq}
\end{multline*}
Note that $n_1n_2=\square$ if and only if $n_1=n_2$ in the Rademacher case, thus we have that the last sum is equal to \begin{align*}
    -2\sum_{1\leq l_i\leq k} \frac{c_{l_1}}{\sqrt{\lvert \mathcal A_{N_{l_1}}\rvert}}\frac{c_{l_2}}{\sqrt{\lvert \mathcal A_{N_{l_2}}\rvert}} \min\{\lvert \mathcal A_{N_{l_1}}\rvert, \lvert \mathcal A_{N_{l_2}}\rvert \} 
     & = -2  + O\left( \sum_{1\leq l_1 < l_2\leq k} \lvert c_{l_1}\overline{c}_{l_2}\rvert\sqrt{\frac{\lvert \mathcal A_{N_{l_1}}\rvert}{\lvert \mathcal A_{N_{l_2}}\rvert}}\right)  
      \\ & = -2 + O\left(k\max_{1\leq l_1<l_2\leq k} \sqrt{\frac{\lvert \mathcal A_{N_{l_1}}\rvert}{\lvert \mathcal A_{N_{l_2}}\rvert}}\right)
\end{align*} by using condition (III). (Again, the error $k\max_{1\leq l_1<l_2\leq k} \sqrt{\frac{\lvert \mathcal A_{N_{l_1}}\rvert}{\lvert \mathcal A_{N_{l_2}}\rvert}}$ is not present if (III$_b$) is satisfied.) Therefore 
\begin{multline*}
    B = -1+ \sum_{p, q} \sum_{1\leq l_i\leq k} \frac{c_{l_1}}{\sqrt{\lvert \mathcal A_{N_{l_1}}\rvert}}\cdots \frac{{c}_{l_4}}{\sqrt{\lvert \mathcal A_{N_{l_4}}\rvert}}\sum_{\substack{n_i\in\mathcal A_{N_{l_i}} \\ n_1n_2n_3n_4 = \square \text{ }(\text{or }n_1n_4=n_2n_3)\\ P^+(n_1) = P^+(n_2)=p \\ P^+(n_3) = P^+(n_4) =q}} 1 \\ +  O\left(k\max_{1\leq l_1<l_2\leq k} \sqrt{\frac{\lvert \mathcal A_{N_{l_1}}\rvert}{\lvert \mathcal A_{N_{l_2}}\rvert}}\right).
\end{multline*}
In the above sum, note that $l_1=l_2, l_3 = l_4$ with the corresponding trivial solutions (after moving the sum over $p, q$ inside) contributes $1$. Moreover, the contribution from non-trivial solutions using (I) is 
\begin{equation*}
\ll\varepsilon_1 \sum_{1\leq l_i\leq k} \lvert c_{l_1}\rvert\cdots \lvert {c}_{l_4}\rvert \ll \varepsilon_1 k^2.
\end{equation*} Hence 
\begin{equation*}
     B \ll \varepsilon_1k^2 + k\max_{1\leq l_1<l_2\leq k} \sqrt{\frac{\lvert \mathcal A_{N_{l_1}}\rvert}{\lvert \mathcal A_{N_{l_2}}\rvert}} +  \sum_{\substack{1\leq l_i\leq k \\ l_1\neq l_2 \text{ or } l_3\neq l_4}} \frac{\lvert c_{l_1}\rvert}{\sqrt{\lvert \mathcal A_{N_{l_1}}\rvert}}\cdots \frac{\lvert {c}_{l_4}\rvert}{\sqrt{\lvert \mathcal A_{N_{l_4}}\rvert}}\sum_{\substack{n_i\in\mathcal A_{N_{l_i}} \\ n_1=n_3, n_2 = n_4 \\ P^+(n_1) = P^+(n_2) }} 1.
\end{equation*}
As before, $n_1=n_3$ implies $n_1\in\mathcal A_{N_{l_1}}\cap \mathcal A_{N_{l_3}}$ and similarly for $n_2=n_4$, thus we employ condition (II) and Cauchy-Schwarz to bound the sum by $\varepsilon_2k^2$ (or just $\varepsilon_2$ if (III$_b$) is satisfied). We conclude that
\begin{equation}\label{eq:Bbound}
B \ll k^2(\varepsilon_1 + \varepsilon_2) + k \max_{1\leq l_1<l_2\leq k} \sqrt{\frac{\lvert \mathcal A_{N_{l_1}}\rvert}{\lvert \mathcal A_{N_{l_2}}\rvert}},
\end{equation}
or 
\begin{equation}\label{eq:BboundDisjoint}
B \ll  k^2\varepsilon_1  + \varepsilon_2
\end{equation}
if (III$_b$) is satisfied.

Combining \eqref{eq:Abound}, \eqref{eq:AboundDisjoint}, \eqref{eq:Bbound}, \eqref{eq:BboundDisjoint} and using Proposition \ref{prop:quantCLT},
we obtain the following approximate joint Gaussianity result. 

\begin{theorem}\label{thm:approxJointGaussian}
    Let $f$ be a Rademacher or Steinhaus multiplicative function, and suppose the $\mathcal A_{N_l}$ for $1\leq l\leq k$  satisfy conditions (I), (I'), (II) and (III), with the values $\varepsilon_1, \varepsilon_1', \varepsilon_2$, respectively. Then, for any real coefficients $c_1, c_2, \ldots, c_k$ with $\sum_{l=1}^k \lvert c_l\rvert^2 = 1$, we have
    \begin{multline*}
\left\lvert\mathbb P\left(\sum_{l=1}^k c_l S_{N_{l}}\leq x \right) - \frac{1}{\sqrt{2\pi}}\int_{-\infty}^x e^{-\frac{t^2}{2}}\emph{d}t\right\rvert \\ \ll \frac{1}{1+\lvert x\rvert^{\frac{16}{5}}}\left(   \varepsilon_2 + k^2(\varepsilon_1 + \varepsilon_1' ) + (1-1_{\text{(III$_b$)}})\left(k^2\varepsilon_2 + k \max_{1\leq l_1<l_2\leq k} \sqrt{\frac{\lvert \mathcal A_{N_{l_1}}\rvert}{\lvert \mathcal A_{N_{l_2}}\rvert}}\right) \right)^{\frac{1}{5}},
 \end{multline*}
 where $1_{\text{(III$_b$)}}$ denotes the indicator function of (III$_b$).
\end{theorem}

In certain cases, it will be necessary to have an analogous result with $B=0$, that is, an error independent of $\varepsilon_1$. For this, we introduce an additional assumption:
\begin{itemize}
    \item[(IV)] Let $P$ be a parameter and suppose that the sets $\mathcal A_{N_l}$ ($1\leq l\leq k$) consist only of integers divisible by some unique prime greater than $P$.
\end{itemize}
Under the assumption (IV), we condition on all the primes $\leq P$ and note that 
\begin{equation*}
    \frac{1}{V_c}\sum_{p>P}\sum_{l=1}^k c_l M_{p, N_l}
\end{equation*} is a martingale (in fact it is a weighted sum of the independent random variables $f(p)$ with $p>P$), where \begin{equation}\label{eq:V_c}
    V_c =\sqrt{\widetilde{\mathbb E}\left(\sum_{p>P}\sum_{l=1}^k c_l M_{p, N_l} \right)^2 } = \sqrt{\sum_{l_1\leq l_2}c_{l_1}c_{l_2}\widetilde{\mathbb E}S_{N_{l_1}}S_{N_{l_2}}}.
\end{equation} Here and throughout $\widetilde{\mathbb E}$ denotes conditional expectation being taken only over primes $>P$. 

\begin{theorem}\label{thm: complicatedapproxJointG}
    Let $f$ be a Rademacher (or Steinhaus) random multiplicative function, and suppose the $\mathcal A_{N_l}$ for $1\leq l\leq k$ satisfy conditions (I'), (II), (III), and (IV), with the values $\varepsilon_1', \varepsilon_2$, and $P$, respectively. With probability \begin{equation}\label{eq:bigProb}1 - O\left(k\sqrt{(\varepsilon_1' + \varepsilon_2)} +(1-1_{\text{(III$_b$)}}) k \max_{1\leq l_1<l_2\leq k} \sqrt{\frac{\lvert \mathcal A_{N_{l_1}}\rvert}{\lvert \mathcal A_{N_{l_2}}\rvert}}\left(k\sqrt{(\varepsilon_1' + \varepsilon_2)}\right)^{1/12}\right)\end{equation} over the $f(p)$ for primes $p\leq P$, we have 
    \begin{equation*}\left\lvert\widetilde{\mathbb P}\left(\sum_{p>P}\sum_{l=1}^k c_l M_{p, N_l}\leq x \right) - \frac{1}{\sqrt{2\pi V_c^2}}\int_{-\infty}^x e^{-\frac{t^2}{2V_c^2}}\emph{d}t\right \rvert  \ll \frac{\left(k\sqrt{\varepsilon_1' + \varepsilon_2}\right)^{\frac{1}{10}}}{1+\lvert x\rvert^{\frac{16}{5}}}, \end{equation*} where $\widetilde{\mathbb P}$ denotes conditional probability being taken only over $f(p)$ for primes $p>P$. 
\end{theorem}

\begin{proof}
    We begin by applying Proposition \ref{prop:quantCLT} to the martingale 
    \begin{equation*}
         \frac{1}{V_c}\sum_{p>P}\sum_{l=1}^k c_l M_{p, N_l}.
    \end{equation*} Noting that  
    \begin{equation*}
        \frac{1}{V_c^2}\widetilde{\mathbb E}\sum_{p>P}\left( \sum_{l=1}^k c_l M_{p, N_l}\right)^2-1=0
    \end{equation*} by the definition of $V_c$, after applying the change of variables $V_cx\mapsto x$, we have
    \begin{multline*}
        \left\lvert\widetilde{\mathbb P}\left(\sum_{p>P}\sum_{l=1}^k c_l M_{p, N_l}\leq x \right) - \frac{1}{\sqrt{2\pi V_c^2}}\int_{-\infty}^x e^{-\frac{t^2}{2V_c^2}}\emph{d}t\right \rvert \\ \ll \frac{\left(V_c^{12}\cdot \widetilde{\mathbb E}\sum_{p>P}\left\lvert \sum_{l=1}^k c_l M_{p, N_l}\right\rvert^4\right)^{\frac{1}{5}}}{1+\lvert x\rvert^{\frac{16}{5}}}.
    \end{multline*} 
    %It is worth noting here that the quantities on the right-hand side are random, but deterministic upon conditioning on the $f(p)$ for $p\leq P$. 
    Now, by exactly the same computations leading to Theorem \ref{thm:approxJointGaussian}, we have 
    \begin{equation*}
        \mathbb E\widetilde{\mathbb E}\sum_{p>P}\left\lvert \sum_{l=1}^k c_l M_{p, N_l}\right\rvert^4\ll k^2(\varepsilon_1' + \varepsilon_2),
    \end{equation*} and 
    \begin{equation*}
        \mathbb EV_c^2 \ll 1 + (1-1_{\text{(III$_b$)}})k \max_{1\leq l_1<l_2\leq k} \sqrt{\frac{\lvert \mathcal A_{N_{l_1}}\rvert}{\lvert \mathcal A_{N_{l_2}}\rvert}}.
    \end{equation*}
        It follows by Markov's inequality with the first moment that with probability as in \eqref{eq:bigProb}, both 
        \begin{equation*}
            \widetilde{\mathbb E}\sum_{p>P}\left\lvert \sum_{l=1}^k c_l M_{p, N_l}\right\rvert^4 \leq k\sqrt{\varepsilon_1'+\varepsilon_2}\end{equation*} and \begin{equation*}
            V_c^2\ll \left(k\sqrt{\varepsilon_1' + \varepsilon_2}\right)^{-1/12}
        \end{equation*} hold, thus finishing the proof. 
\end{proof}

The advantage of introducing condition (IV) and Theorem \ref{thm: complicatedapproxJointG} is that the parameters $\varepsilon_1'$ and $\varepsilon_2$ can be taken to be much smaller than before (depending on the size of $P$). This leads to a very precise approximation of the conditional probabilities over primes bigger than $P$, at least over most realizations of the $f(p)$ for $p\leq P$. 

\subsection{Multivariate Central Limit Theorem} Finally, we use Proposition \ref{prop:multivarCLT} to bootstrap Theorems \ref{thm:approxJointGaussian} and \ref{thm: complicatedapproxJointG} to quantitative multivariate central limit theorems.     

\begin{theorem}\label{thm:MCLT} Let $f$ be a Rademacher (or Steinhaus) random multiplicative function, and suppose the $\mathcal A_{N_l}$ for $1\leq l\leq k$ satisfy conditions (I), (I'), (II) and (III), with the values $\varepsilon_1, \varepsilon_1', \varepsilon_2$, respectively. Further, let $X=(S_{N_1}, S_{N_2}, \ldots, S_{N_k})$ (recall \eqref{eq:S_Ndef} for the definition of $S_N$) and let $Y$ be a standard multivariate Gaussian random vector with independent coordinates. Then 
\begin{equation}
 \sup_{\|u\|_{\emph{Lip}}\leq 1} \lvert \mathbb EuX - \mathbb EuY\rvert\ll k^{\frac{1}{2}}  \left(k^2\varepsilon_2 + k \max_{1\leq l_1<l_2\leq k} \sqrt{\frac{\lvert \mathcal A_{N_{l_1}}\rvert}{\lvert \mathcal A_{N_{l_2}}\rvert}} \right)^{\frac1{5(k+1)}}.
 \end{equation}
If the $\mathcal A_{N_l}$ satisfy (I'), (II), (III), and (IV), then with probability as in \eqref{eq:bigProb} over the $f(p)$ for primes $p\leq P$, we have 
 \begin{equation}
 \sup_{\|u\|_{\emph{Lip}}\leq 1} \lvert \widetilde{\mathbb E}uX - \mathbb EuY'\rvert \\ \ll k^{\frac{1}{2}} \left(k\sqrt{\varepsilon_1' + \varepsilon_2}\right)^{\frac1{10(k+1)}},
 \end{equation}
 where $Y'$ is a multivariate Gaussian random vector with the covariance structure $\mathbb EY'_iY'_j = \widetilde{\mathbb E}S_{N_{l_i}} S_{N_{l_j}}.$
\end{theorem}

\begin{proof}
This follows directly from Proposition \ref{prop:multivarCLT} with $p=q=2$, using the conclusions of Theorems \ref{thm:approxJointGaussian} or \ref{thm: complicatedapproxJointG} to bound the right-hand side. Moreover, note that both norms are proportional to $\sqrt k$ (in the conditional case, this happens with probability as in \eqref{eq:bigProb} by exactly the same computation as in the proof of Theorem \ref{thm: complicatedapproxJointG}). It may be helpful to observe that 
\begin{equation*}
    \mathbb P \left(\sum_{l=1}^k c_lY_l'\leq x \right) =\frac{1}{\sqrt{2\pi V_c^2}} \int_{-\infty}^x e^{-\frac{t^2}{2V_c^2}}\text{d}t,
\end{equation*}
for by definition $\sum_{l=1}^k c_lY_l'$ is a Gaussian random variable with mean $0$ and variance $V_c^2$ (recall \eqref{eq:V_c}). 
\end{proof}

\begin{remark}
    As we remarked before, assumption (IV) places us in a position where 
    \begin{equation*}
         \frac{1}{V_c}\sum_{p>P}\sum_{l=1}^k c_l M_{p, N_l}
    \end{equation*}
    is conditionally a weighted sum of independent random variables. Therefore, instead of proving Theorem \ref{thm: complicatedapproxJointG} and bootstrapping it to Theorem \ref{thm:MCLT}, we could have opted to use the Normal Approximation Result 1 from the work of Harper \cite{Harper2023_ASLF}, which yields a stronger result in terms of the dependence on $k$. However, our argument only utilizes the martingale structure and may potentially be useful in other situations with a weaker assumption than (IV). 
\end{remark}

Theorem \ref{thm:MCLT} may then be used to estimate the probability that $\max_{1\leq l\leq k} S_{N_l}\leq t$ in terms of the probability that $\max_{1\leq l\leq k} Y_{l}\leq t + O(1)$. 
\begin{cor}\label{cor:normalcomparison}
    Adapt the notation and conditions in the statement of Theorem \ref{thm:MCLT}. Then 
    \begin{multline*}
        \mathbb P\left(\max_{1\leq l\leq k} S_{N_l} \leq t\right)\leq \mathbb P\left(\max_{1\leq l\leq k} Y_l\leq t+\eta\right) \\ + O\left(k^{\frac{3}{2}}\eta^{-1} \left(  k^2(\varepsilon_1 + \varepsilon_1'  + \varepsilon_2) + k \max_{1\leq l_1<l_2\leq k} \sqrt{\frac{\lvert \mathcal A_{N_{l_1}}\rvert}{\lvert \mathcal A_{N_{l_2}}\rvert}} \right)^{\frac1{5(k+1)}}\right).
    \end{multline*}
    If additionally (IV) is satisfied (but not necessarily (I)), then with probability as in \eqref{eq:bigProb} over the $f(p)$ for primes $p\leq P$, we have 
    \begin{equation*}
        \widetilde{\mathbb P}\left(\max_{1\leq l\leq k} S_{N_l} \leq t\right)\leq \mathbb P\left(\max_{1\leq l\leq k} Y_l'\leq t+\eta\right)  + O\left(k^{\frac{3}{2}}\eta^{-1} \left(k\sqrt{\varepsilon_1' + \varepsilon_2}\right)^{\frac1{10(k+1)}}\right).
    \end{equation*}
    
\end{cor}

\begin{proof}
    Take $u(x_1, \ldots, x_k) = \prod_{i=l}^k s(x_l)$ in Theorem \ref{thm:MCLT}, where $s$ is a smooth function\footnote{One possibility is $s(x) = \phi(t+\eta - x)/(\phi(t+\eta - x) + \phi(x-t))$ with $\phi(x)=\exp(-1/x)$ for $x>0$ and $\phi(x)=0$ for $x\leq 0$.} such that $s(x) = 1$ if $x\leq t$ and $s(x) = 0$ if $x>t+\eta$. Note that the Lipschitz constant of $u$ is bounded by $k\eta^{-1}$, so after rescaling $u$ by $k\eta^{-1}$, we get an extra $k\eta^{-1}$ in the error. Observe that the difference between the desired probabilities is bounded by $\mathbb E u(S_{N_1}, \ldots, S_{N_k}) - \mathbb EuY$, thus finishing the proof.
\end{proof}


\subsection{Slow variation} Here we shall prove a general slow variation property of sums of random multiplicative functions, which is also applicable in short intervals. This generalizes a result found in the work of Lau, Tenenbaum and Wu \cite{LauTenenbaumWu2013MeanValuesRMF}, making use of Lemma \ref{lemma:billings}. In this paper, we will use this proposition in proving the almost sure upper bound in the case of polynomial images, but it is relevant in proving almost sure upper bounds in other cases too. \begin{prop}\label{prop:slowVar}
Let $f$ be a Rademacher (or Steinhaus) random multiplicative function. Let $\mathcal A$ be an infinite subset of the square-free (or all) positive integers, and suppose that $\mathcal A_N=\mathcal A\cap (N-H, N]$ for some $N\geq H=H(N)\to\infty$ as $N\to\infty$. For any increasing sequence $(N_l)_{l=1}^\infty$ such that $N_1$ is large enough and $N_{l+1}-H_{l+1}\leq N_l$ for all $l$, we have
\begin{multline*}
\mathbb P\left(\max_{N_l<N\leq N_{l+1}}\left\lvert\sum_{n\in\mathcal A_N\setminus\mathcal A_{N_l}}f(n)\right\rvert\geq\lambda\right) \\  \ll \frac{1}{\lambda^4}\left(\left(\sum_{\substack{N_l - H_l<n\leq N_{l+1} - H_{l+1} \\ n\in\mathcal A}} \tau_3(n)\right)^2 + \left(\sum_{\substack{N_l<n\leq N_{l+1} \\ n\in\mathcal A}} \tau_3(n)\right)^2 \right).
\end{multline*} 
\end{prop}

\begin{proof}
Suppose that 
\begin{equation*}
    \max_{N_l<N\leq N_{l+1}}\left\lvert\sum_{n\in\mathcal A_N\setminus\mathcal A_{N_l}}f(n)\right\rvert\geq \lambda.
\end{equation*} 
Writing 
\begin{equation*}
    \sum_{n\in\mathcal A_N\setminus\mathcal A_{N_l}}f(n) = -\sum_{\substack{N_l-H_l<n\leq N-H \\ n\in\mathcal A}}f(n) + \sum_{\substack{N_l<n\leq N \\ n\in\mathcal A}}f(n),
\end{equation*} it follows by the triangle inequality that 
\begin{equation*}
    \max_{N_l<N\leq N_{l+1}}\left\lvert \sum_{\substack{N_l-H_l<n\leq N-H \\ n\in\mathcal A}}f(n)\right\rvert \geq \frac{\lambda}2 \hspace{0.5cm}\text{ or }\hspace{0.5cm} \max_{N_l<N\leq N_{l+1}}\left\lvert \sum_{\substack{N_l<n\leq N \\ n\in\mathcal A}}f(n)\right\rvert\geq \frac{\lambda }{2}.
\end{equation*} 
Let $M_l = N_l - H_l$. Note that for every $u<v$ we have
\begin{equation*}
    \mathbb P\left(\left\lvert \sum_{\substack{u < n\leq v \\ n\in\mathcal A}}f(n) \right\rvert\geq\lambda\right)\leq \frac{1}{\lambda^4}\left(\sum_{\substack{u<n\leq v \\ n\in\mathcal A}} \tau_3(n)\right)^2
\end{equation*} by the rough hypercontractive inequality (see \cite[Probability Result 1]{Harper2019_MomentsII} for example). Employing Lemma \ref{lemma:billings} for $\sum_{M_l<n\leq M_{l+1}}f(n)$ and $\sum_{N_l<n\leq N_{l+1}}f(n)$ gives 
\begin{multline*}
    \mathbb P\left(\max_{N_l<N\leq N_{l+1}}\left\lvert\sum_{n\in\mathcal A_N\setminus\mathcal A_{N_l}}f(n)\right\rvert\geq \lambda\right) \\ \ll \frac{1}{\lambda^4}\left(\left(\sum_{\substack{M_l<n\leq M_{l+1} \\ n\in\mathcal A}} \tau_3(n)\right)^2 + \left(\sum_{\substack{N_l<n\leq N_{l+1} \\ n\in\mathcal A}} \tau_3(n)\right)^2 \right),
\end{multline*}
finishing the proof. 
\end{proof}


\section{Proof of Theorem \ref{thm:asboundsPoly}: Polynomial images}
To prove our theorem for polynomial images, we will only need Theorem \ref{thm:approxJointGaussian} and the simpler part of Corollary \ref{cor:normalcomparison}. Let $P$ be a polynomial as in the statement of Theorem \ref{thm:asboundsPoly} and consider\footnote{We are instead considering $\mathcal A_N\subseteq [1, c_PN^{\deg P}]$ for convenience due to the scaling of polynomial images.} $\mathcal A_N = \{P(n) : n\leq N, P(n)>0\}$. Note that $\lvert \mathcal A_N\rvert \sim \kappa_PN$ for some $\kappa_P>0$ depending\footnote{In the Steinhaus case we have $\kappa_P=1$, whereas in the Rademacher case, $\kappa_P$ is the density of square-free values of $P$.} on $P$. Recall 
\begin{equation}\label{eq:S_NdefPoly}
     S_N = \frac{1}{\sqrt{\lvert \mathcal A_N\rvert}}\sum_{n\in\mathcal A_N} \mathfrak Tf(n) \sim \frac{1}{\sqrt{\kappa_P N}} \sum_{\substack{n\leq N}} \mathfrak Tf(P(n)).
\end{equation}
Moreover, for a large parameter $X$, let $N_l = \lambda^l X$ with $\lambda = \exp\left(\sqrt{\log X}\right)$, where $1\leq l\leq k = (\log X)^{\varepsilon_0}$ with $\varepsilon_0>0$ to be specified later. We begin by checking that these sets satisfy properties (I), (I'), (II) and (III). 

\subsubsection*{(I)} We employ the result of \cite{ChinisShala2025}
 (\cite{KlurmanShkredovXu2023}), which we recall here. 
 
 \begin{prop}
 Let $P$ be a polynomial as in the statement of Theorem \ref{thm:asboundsPoly}. There exists a constant $\delta_P>0$ such that the number of non-trivial solutions to the equation 
 \begin{equation*}
     P(n_1)P(n_2)P(n_3)P(n_4) = \square \hspace{0.5cm} (\text{or }P(n_1)P(n_2) = P(n_3)P(n_4))
 \end{equation*} with $n_i\leq N$ is $\ll N^{2-\delta_P}$. 
 \end{prop}
 
Using the above with $N = \max\{N_{l_1}, \ldots, N_{l_4}\}\ll \lambda^k X$ allows us to take $\varepsilon_1 \ll X^{-\delta_P'}$ for some slightly smaller fixed $\delta_P'>0$ (possibly depending on $P$), since $\lambda$ is so that $\lambda^k\ll_{\epsilon}X^{\epsilon}$ for all $\epsilon>0$. 

\subsubsection*{(I')} We will simply take $\varepsilon_1' =\varepsilon_1$.  

\subsubsection*{(II)} 
    We check condition (II) as follows. We introduce a smoothness parameter $y$, and split the range of $P^+(n_1) = P^+(n_2)=p$ to $p\leq y$ and $p>y$. The latter range also contains very large primes, which are dealt with separately. In any case, we show the contribution of primes $p>y$ is bounded above by 
    \begin{equation*}
\sum_{p>y}\frac{\lvert\mathcal A_{N_{l_1}}\rvert\lvert\mathcal A_{N_{l_2}}\rvert}{p^2}\ll \frac{\lvert \mathcal A_{N_{l_1}}\rvert\lvert\mathcal A_{N_{l_2}}\rvert}{y}.
    \end{equation*}
    The contribution of primes $p\leq y$ is bounded by the product of the number of $y$-smooth values of $n\in\mathcal A_N$ and the number of $y$-smooth values of $m\in\mathcal A_M$. Optimizing the choice of $y$ yields the saving $\varepsilon_2$. 
    
Concretely, using the calculations from the end of \cite[Section 4]{ChinisShala2025} (\cite[Section 2]{KlurmanShkredovXu2023}), we may take 
$$\varepsilon_2 \ll \frac{1}{y} + \frac{\psi_P(N_{l_1}, y)\psi_P(N_{l_2}, y)}{N_{l_1}N_{l_2}},$$ where $\psi_P(N, y)$ is the number of $n\leq N$ such that $P(n)$ is $y$-smooth. By a result of Hmyrova \cite{Hmyrova1966}, when $y\geq\log N$ we have the bound\footnote{The bound there is stated for irreducible $P$, but obviously the number of values of $n\leq N$ for which $P(n)$ is $y$-smooth is bounded above by the number of $n\leq N$ for which some irreducible factor of $P(n)$ is $y$-smooth.} $\psi_P(N, y)\ll N(e/u)^u$, where $u=\log N/\log y$. After an unpleasant exercise in manipulating logarithms and optimizing in $y$, 
we choose $y  = \exp\left(\sqrt{\log X\log\log X}\right)$. This gives 
$\varepsilon_2 \ll \exp(-\sqrt{\log X\log\log X})$. 

\subsubsection*{(III)} Note that the sets $\mathcal A_N$ are nested, so (III$_a$) is satisfied. 

\subsection{Almost sure upper bound} We now have everything we need to prove the almost sure upper bound. We first apply Theorem \ref{thm:approxJointGaussian} with $k=1$ (at a single scale), yielding the following quantitative central limit theorem. 

\begin{cor}\label{cor:quantCLTpoly}
    Let $f$ be a Rademacher (or Steinhaus) random multiplicative function, and let $P\in\mathbb Z[x]$ be a polynomial as in Theorem \ref{thm:asboundsPoly}. We have
     \begin{multline*}
\left\lvert\mathbb P\left(\frac{1}{\sqrt{\kappa_P N}}\sum_{n\leq N}\mathfrak T f(P(n))\leq x \right) - \frac{1}{\sqrt{2\pi}}\int_{-\infty}^x e^{-\frac{t^2}{2}}\emph{d}t\right\rvert \\ \ll \frac{1}{1+\lvert x\rvert^{\frac{16}{5}}}\exp\left(-\frac{1}{5} \sqrt{\log N\log\log N}\right).
    \end{multline*}
\end{cor}

We will use Corollary \ref{cor:quantCLTpoly} at certain test points $(N_l)_{l=1}^\infty$, combined with the following almost sure slow variation result between the test points.

\begin{lemma}\label{lemma:slowVarPoly} Let $f$ be a Rademacher (or Steinhaus) random multiplicative function, and let $P\in\mathbb Z[x]$ be a polynomial as in Theorem \ref{thm:asboundsPoly}. 
For any constant $A>0$, there exists a constant $c>0$ such that for $N_l = \lfloor e^{l^c}\rfloor$, we have that
    \begin{equation*}
        \max_{N_{l}<N\leq N_{l+1}} \left\lvert\sum_{n\leq N} f(P(n)) - \sum_{n\leq N_{l}} f(P(n)) \right\rvert \ll \frac{\sqrt{N_{l+1}}}{(\log N_{l+1})^A}
    \end{equation*} holds almost surely for every $l\in\mathbb N$. 
\end{lemma}

\begin{proof}
    By Proposition \ref{prop:slowVar}, we have 
\begin{equation*}
\mathbb P\left(\max_{N_l<N\leq N_{l+1}}\left\lvert\sum_{N_l<N\leq N}f(P(n))\right\rvert\geq x\right) \ll \frac{1}{x^4} \left(\sum_{\substack{N_l<n\leq N_{l+1} }} \tau_3(P(n))\right)^2 
\end{equation*}
Lemma \ref{lemma:PolyBound} gives 
\begin{equation*}
        \sum_{N_l<n\leq N_{l+1}} \tau_{3}(P(n)) \ll (N_{l+1} - N_l)^2(\log N_{l+1})^{c_P},
    \end{equation*} 
    for some constant $c_P>0$ depending on $P$. Taking $x = \sqrt{N_{l+1}}/(\log N_{l+1})^A$, we obtain 
    \begin{equation*}
        \mathbb P\left(\max_{N_l<N\leq N_{l+1}}\left\lvert\sum_{N_l<N\leq N}f(P(n))\right\rvert\geq x\right) \ll \left(\frac{N_{l+1} - N_l}{N_{l+1}} \right)^2 (\log N_{l+1})^{c_P +4A}.
    \end{equation*}
Plugging in $N_l = \lfloor e^{l^c}\rfloor$ and summing over $l$ to get a convergent series, we see that choosing $0<c<1/({c_P+4A+2})$ suffices. Applying Borel-Cantelli finishes the proof.
\end{proof} 

We finally prove the almost sure upper bound. By Lemma \ref{lemma:slowVarPoly} with $A=1$, it suffices to show that almost surely there exists a constant $C>0$ such that for all $l\in\mathbb N$ we have
    \begin{equation}\label{eq:ubPoly}
        \left\lvert\sum_{n\leq N_l} f(P(n))\right\rvert\leq C \sqrt{N_l\log\log N_l},
    \end{equation}
    where $N_l$ is as in the statement of the lemma (in the Steinhaus case, we do this separately for the real and imaginary parts). By Corollary \ref{cor:quantCLTpoly}, the probability that the above fails for a fixed $l$ and large enough implied constant $C$ is $\ll (\log N_l)^{-\frac{C^2}{2}}$. Making $C$ larger (in terms of $c$) if necessary, we have that the series of probabilities $$\sum_{l=1}^\infty (\log N_l)^{-\frac{C^2}{2}} = \sum_{l=1}^\infty \frac{1}{l^{\frac{cC^2}{2}}}<\infty,$$ thus \eqref{eq:ubPoly} holds almost surely by Borel-Cantelli (with a possibly even larger constant $C$ to account for the finitely many failures). 



    

\subsection{Almost sure lower bound} Here we will use Corollary \ref{cor:normalcomparison} with the scales $N_1, N_2, \ldots, N_k$ defined at the start of the section. In fact, we will prove something slightly stronger.

\begin{theorem}\label{thm:localizedBoundPoly}
    Let $f$ be a Rademacher (or Steinhaus) random multiplicative function and let $P\in\mathbb Z[x]$  be a polynomial as in Theorem \ref{thm:asboundsPoly}. There exists a constant $c>0$ such that for any large enough $X$, we have that 
    \begin{equation*}
        \max_{X<N\leq X^2} \frac{1}{\sqrt N} \left\lvert\sum_{n\leq N}f(P(n)) \right\rvert \gg \sqrt{\log\log X}
    \end{equation*} holds with a suitably small implied constant with probability $1-O(\exp(-(\log X)^{c}))$. 
\end{theorem}

In light of Lemma \ref{lemma:asLbGeneral}, this suffices to establish the almost sure lower bound in Theorem \ref{thm:asboundsPoly}. Henceforth we prove Theorem \ref{thm:localizedBoundPoly}. Applying Corollary \ref{cor:normalcomparison} we obtain the following normal comparison result.

\begin{cor}\label{cor:MCLTpoly}
    Let $X=(S_{N_1}, S_{N_2}, \ldots, S_{N_k})$ (recall \eqref{eq:S_NdefPoly} for the definition of $S_N$), and let $Y$ be a standard multivariate Gaussian random vector with independent coordinates. Then 
\begin{equation*}
\mathbb P\left(\max_{1\leq l\leq k} S_{N_l} \leq t\right)\leq \mathbb P\left(\max_{1\leq l\leq k} Y_l\leq t+\eta\right) + O\left( \eta^{-1}\exp\left(-\frac{1}{11}(\log X)^{\frac{1}{3}} \right)\right).
 \end{equation*}
\end{cor}
\begin{proof}
    Recalling that $k=(\log X)^{\varepsilon_0}$ and $\lambda = \exp(\sqrt{\log X})$, Corollary \ref{cor:normalcomparison} furnishes the bound 
    \begin{multline*}
        \frac{(\log X)^{\frac{3\varepsilon_0}{2}}} {\eta}\Big((\log X)^{2\varepsilon_0}\exp\left(-\frac{1}{5}\sqrt{\log X\log\log X} \right)  + X^{-\delta_P}(\log X)^{2\varepsilon_0} \\ + (\log X)^{\varepsilon_0}\exp\left(-\frac{1}{2}\sqrt{\log X}\right)  \Big)^{\frac1{5(\log X)^{\varepsilon_0}}},
    \end{multline*}
    which upon inspection is less than the claimed bound if $\varepsilon_0>0$ is sufficiently small. 
\end{proof}

Now the probability $\mathbb P\left(\max_{1\leq l\leq k} Y_l\ll\sqrt{\log k}\right)$ can be suitably estimated using Lemma \ref{lemma:normalsBigLemma}. Recall that $\mathbb E Y_{l_1}Y_{l_2}= 0$ for $i\neq j$ since the components of $Y$ are independent, so we have 
\begin{equation*}
    \mathbb P\left(\max_{1\leq l\leq k} Y_l\leq\sqrt{(2-\delta)\log k}\right)\leq \exp\left(-\Theta\left(\frac{k^{\delta/20}}{\sqrt{\log k}}\right)\right).
\end{equation*}
Recalling that $k=(\log X)^{\varepsilon_0}$ and taking $\eta>0$ and $\delta>0$ fixed, we obtain 
\begin{equation*}
    \mathbb P\left(\max_{1\leq i\leq k} S_{N_l} \gg \sqrt{\log\log X}\right)  \geq 1 - O\left(\exp\left(-(\log X)^{O(\varepsilon_0)}\right)\right)
\end{equation*}
for a suitably small implied constant (making $\varepsilon_0>0$ smaller if necessary). Theorem \ref{thm:localizedBoundPoly} follows by noticing that $X\leq \lambda X\leq N_l\leq \lambda^kX\leq \exp\left((\log X)^{\frac{1}{2}+\varepsilon_0}\right)X\leq X^2$ for all $1\leq l\leq k$. 


\section{Proof of Theorem \ref{thm:asBoundsShort}: Short intervals} Here we will require the more sophisticated part of Corollary \ref{cor:normalcomparison}. Let $H$ be as in the statement of Theorem \ref{thm:asBoundsShort} and consider $\mathcal A_N = \{n\in\mathbb N : N-H<n\leq N\}$. Note that $\lvert\mathcal A_N\rvert\sim \kappa H$ for some\footnote{In the Steinhaus case we have $\kappa = 1$, whereas in the Rademacher case $\kappa = \frac{6}{\pi^2}$.} $\kappa>0$. Recall 
\begin{equation*}
    S_N = \frac{1}{\sqrt{\lvert\mathcal A_N\rvert}}\sum_{n\in\mathcal A_N} \mathfrak Tf(n)\sim \frac{1}{\sqrt{\kappa H}}\sum_{N-H<n\leq n }\mathfrak Tf(n).
\end{equation*}
Moreover, let $X$ be a large parameter. For some small $\delta>0$ depending on the function $H$ and then even smaller $\varepsilon_0>0$, find $k= (X/H(X))^{\varepsilon_0}$ primes $l$ in the range $h(X)/2<l\leq h(X)$, where $h(X) := (X/H(X))^{\delta}$. We will not fix $\delta$ or $\varepsilon_0$ for now, as there will be various (but finitely many) places where we may have make them smaller. For these values of $l$, let $N_l = lX$ with corresponding $H_l = H(N_l)$. 
Observe that for large enough $X$ and small enough $\delta>0$, the sets $\mathcal A_{N_l}$ are disjoint. It is not these sets that we will apply Theorem \ref{thm: complicatedapproxJointG} to. If we did this, the very smooth integers in the sets $\mathcal A_{N_l}$ would prevent us from obtaining a suitably small value for $\varepsilon_2$. We bypass this by splitting $S_{N_l}$ as 
\begin{equation}\label{eq:splitting1}
    \sum_{N_l - H_l<n\leq N_l}f(n) = \sum_{(Xh(X))^{\frac{2}{3}}< p}f(p)\sum_{\substack{N_l-H_l < n\leq N_l \\ p\mid n}}f\left(\frac np\right) + \sum_{\substack{N_l-H_l<n\leq N_l \\ P^+(n)\leq (Xh(X))^{\frac{2}{3}}}}f(n).
\end{equation}
If $H_1> N_1/(\log N_1)^2$ (hence since $H$ is increasing we have $H_l \gg Xh(X)/(\log X)^2$ for all values of $l$), we choose some $\varepsilon>0$ to be specified later and additionally split the first part as 
\begin{equation}\label{eq:splitting2}
    \sum_{N_l - H_l<n\leq N_l}f(n) = \sum_{(Xh(X))^{\frac{2}{3}}< p }f(p)\sum_{\substack{N_l-H_l < n\leq N_l \\ p\mid n \\ \Omega(n)\leq (1+\varepsilon)\log\log X}}f\left(\frac np\right) \\ + \sum_{\substack{N_l-H_l<n\leq N_l \\  P^+(n)\leq (Xh(X))^{\frac{2}{3}} \\ \text{or } \Omega(n)> (1+\varepsilon)\log\log X }}f(n).
\end{equation}
We first show that the sum over integers with $\Omega(n)> (1+\varepsilon)\log\log X$ above may be typically ignored for all values of $l$. 

\begin{lemma}\label{lemma:ignoreTooManyPrimesShort}
Assume that $H_1 > N_1/(\log N_1)^2$ and let $\varepsilon>0$. There exists $\varepsilon' = \varepsilon'(\varepsilon)$ such that with probability $1-O((\log X)^{-\varepsilon'+2\varepsilon_0})$ we have 
\begin{equation*}
    \frac{1}{\sqrt{H_l}}\sum_{\substack{N_l-H_l<n\leq N_l \\ \Omega(n)>(1+\varepsilon)\log\log X}} f(n)\ll1
\end{equation*} for all of the $k$ chosen values of $h(X)/2<l\leq h(X).$\end{lemma}

\begin{proof}
    The probability that $$\left\lvert\sum_{\substack{N_l - H_l < n\leq N_l \\ \Omega(n)>(1+\varepsilon)\log\log X}} f(n)\right\rvert\geq\lambda$$ is bounded by $\lambda^{-2}$ times the number of $n\in (N_l - H_l, N_l]$ with $\Omega(n)>(1+\varepsilon)\log\log X$ by Markov's inequality using a second moment estimate. The number of such $n$ is $\ll H_l/(\log X)^{\varepsilon'}$ by Lemma \ref{lemma:tooManyPrimes}. Taking a union bound over $k$ values of  $l$ and recalling that $k=(X/H(X))^{\varepsilon_0}\ll (\log X)^{2\varepsilon_0}$ (making $\varepsilon_0>0$ smaller if necessary) gives a probability $1-O((\log X)^{-\varepsilon'+2\varepsilon_0})$ that all of the sums are bounded. 
\end{proof}

If $H_1\leq N_1/(\log N_1)^2$ (in which case by concavity of the function $H$ we have $H_l$  \linebreak$\ll Xh(X)/(\log X)^2$ for all values of $l$), let $\mathcal B_{N_l}$ be the subset of $\mathcal A_{N_l}$ of those integers $n$ with $P^+(n)>(Xh(X))^{\frac{2}{3}}$. If $H_1> N_1/(\log N_1)^2$, add the restriction $\Omega(n)\leq(1+\varepsilon)\log\log X$. Now we work with 
\begin{equation*}
    \sum_{n\in\mathcal B_{N_l}}f(n) + O\left(\left\lvert\sum_{\substack{N_l-H_l<n\leq N_l \\  P^+(n)\leq (Xh(X))^{\frac{2}{3}}}}f(n)\right\rvert\right).
\end{equation*}
We further show that for many values of $l$, the smooth sum is typically small. 

\begin{lemma}\label{lemma:ignoresmoothShort}
    For large enough $X$, with probability $1-O((\log \log X)^{-\frac{2}{3}})$ over $f(p)$ with $p\leq (Xh(X))^{\frac{2}{3}}$, there exist $\geq \frac{k}{2}$ values of $l$ such that 
    \begin{equation*}
        \frac{1}{\sqrt H_l}\sum_{\substack{N_l-H_l<n\leq N_l \\ P^+(n)\leq (Xh(X))^{\frac{2}{3}}}}f(n)\ll (\log\log X)^{\frac{1}{3}}.    \end{equation*}
\end{lemma} 

\begin{proof}
 The probability that the normalized smooth sum is greater than $\lambda$ is bounded by $1/\lambda^2$ by Markov's inequality using a second moment estimate. We choose  $\lambda = (\log\log X)^{1/3}$ so that the expected number of $l$ with the normalized smooth sum greater than $\lambda$ is $\leq k/\lambda^2$. Then for large enough $X$, with probability $1-O((\log \log X)^{-\frac{2}{3}})$ we may find $\geq \frac{k}{2}$ suitable values of $l$. 
\end{proof}

Henceforth condition on the values of $f(p)$ for $p\leq (Xh(X))^{\frac{2}{3}}$, upon which we may find $\geq \frac{k}{2}$ fixed values of $l$ that we will refer to as \emph{good}, with the property that
\begin{equation*}
    \sum_{N_l - H_l < n\leq N_l}f(n) = \sum_{n\in\mathcal B_{N_l}}f(n) + \sum_{\substack{N_l-H_l<n\leq N_l \\ \Omega(n)>(1+\varepsilon)\log\log X}} f(n) + O\left(\sqrt{H_l}(\log \log X)^{\frac{1}{3}} \right)
\end{equation*} for all good $l$. (The second sum appears only when $H_1 > N_1/(\log N_1)^2$.)



\subsection{Almost sure lower bound}
Throughout this section, all values of $l$ will be good even if we do not state so. As before, we will actually prove something slightly stronger.
\begin{theorem}\label{thm:localizedBoundShort}
    Let $f$ be a Rademacher (or Steinhaus) random multiplicative function and let $H$ be a function as in Theorem \ref{thm:asBoundsShort}. For any large enough $X$, we have that 
    \begin{equation*}
        \max_{X<N\leq X^2} \frac{1}{\sqrt H} \left\lvert\sum_{N-H<n\leq N}f(n) \right\rvert \gg \sqrt{\log \frac{X}{H(X)}}
    \end{equation*} holds with a suitably small implied constant with probability $1-O((\log \log X)^{-\frac{2}{3}})$. 
\end{theorem}

In light of Lemma \ref{lemma:asLbGeneral}, this suffices to establish Theorem \ref{thm:asBoundsShort}. Henceforth we prove Theorem \ref{thm:localizedBoundShort}. 

Let
\begin{equation*}
    \widetilde S_{N_l} = \frac{1}{\sqrt{\lvert \mathcal B_{N_l}\rvert}}\sum_{n\in\mathcal B_{N_l}} \mathfrak Tf(n).
\end{equation*}
We wish to apply Corollary \ref{cor:normalcomparison} to the $\mathcal B_{N_l}$, so we will check conditions (I'), (II), (III), and (IV) for these sets.

\subsubsection*{(I')}  If $n_i\in\mathcal B_{N_{l_i}}$ are such that $n_1n_2n_3n_4 = \square$ (or $n_1n_4 = n_2n_3$) with $P^+(n_1) = \cdots = P^+(n_4)=p > (Xh(X))^{\frac{2}{3}}$, then $n_i' = n_i/p$ satisfy $n_1'n_2'n_
  3'n_4'=\square$ (or $n_1'n_2' = n_3'n_4'$) and $n_i'\leq N_{l_i}/p\leq Xh(X)/p$. For a fixed $p$, the number of such $n_i'$ is $\ll (Xh(X)\log X/p)^2$. Over all primes $p>(Xh(X))^{\frac{2}{3}}$, we have at most 
\begin{equation*}
    \ll (Xh(X)\log X)^2 \sum_{(Xh(X))^{\frac{2}{3}}<p\leq kX}\frac{1}{p^2}\ll (Xh(X))^{\frac{4}{3}}\log X
\end{equation*}
solutions for $n_i$. For  $\varepsilon_1' = X^{-(\frac{2}{3} - \frac{8}{15})}\log X$, this is $\ll \varepsilon_1' H_1^2\ll \varepsilon_1'\sqrt{H_{l_1}\cdots H_{l_4}}$ uniformly in the values of the $l_i$, since $H\geq N^{\frac{11}{15}}.$

\subsubsection*{(II)} The number of $n_1\in\mathcal B_{N_{l_1}}, n_2\in\mathcal B_{N_{l_2}}$ with $P^+m=P^+(n)$ is bounded by 
\begin{multline*}
    \sum_{(Xh(X))^{\frac{2}{3}}<p} \#\{n_1\in\mathcal B_{N_{l_1}} : P^+(n_1) = p \} \cdot \#\{n_2\in\mathcal B_{N_{l_2}} : P^+(n_2) = p \}  \\ \ll \left(\frac{H_{l_1}}{(Xh(X))^{\frac{2}{3}}} + 1\right)\sum_{(Xh(X))^{\frac{2}{3}}<p} \#\{n_2\in\mathcal B_{N_{l_2}} : P^+(n_2) = p \} \\  \ll \left(\frac{H_{l_1}}{(Xh(X))^{\frac{2}{3}}} + 1 \right)H_{l_2}.
\end{multline*}
For $\varepsilon_2 = (Xh(X))^{-\frac{2}{3}}$, this is certainly $\ll \varepsilon_2 H_{l_1}H_{l_2}$ uniformly in the values of $l_1$ and $l_2$. 

\subsubsection*{(III)} As we remarked in the beginning, the stronger assumption (III$_b$) is satisfied here. 

\subsubsection*{(IV)} We have chosen $P = (Xh(X))^{\frac{2}{3}}$ and the rest is true by definition of the $\mathcal B_{N_l}$. 

Now we apply Corollary \ref{cor:normalcomparison} to the $\widetilde{S}_{N_l}$ with good $l$, obtaining that with probability as in \eqref{eq:bigProb} over the primes $p\leq P$ (with the values of $\varepsilon_1', \varepsilon_2$ that we just obtained), we have
  \begin{equation}\label{eq:shortIntsNormal}
        \widetilde{\mathbb P}\left(\max_{l} \widetilde{S}_{N_l} \leq t\right)\leq \mathbb P\left(\max_{l} Y_l\leq t+\eta\right)   + O\left(k^{\frac{3}{2}}\eta^{-1} \left(k\sqrt{\varepsilon_1' + \varepsilon_2}\right)^{\frac1{5(k+1)}}\right),
    \end{equation}
    where the $Y_l$ are Gaussian random variables with the covariance structure 
\begin{equation*}\mathbb E Y_{l_1}Y_{l_2} = \widetilde{\mathbb E}\widetilde{S}_{N_{l_1}}\widetilde{S}_{N_{l_2}} = \frac{1}{\sqrt{\lvert\mathcal B_{N_{l_1}}\rvert\lvert\mathcal B_{N_{l_2}}\rvert}}\sum_{\substack{n_1\in\mathcal B_{N_{l_1}} \\ n_2\in\mathcal B_{N_{l_2}}\\ P^+(n_1) = P^+(n_2)}}\mathfrak Tf(n_1)\mathfrak Tf(n_2).\end{equation*} We would like to apply Lemma \ref{lemma:normalsBigLemma} to the $\widetilde{Y}_l = Y_l/\mathbb EY_l^2$, so we first have to check that there exists a sufficiently small $\epsilon>0$ such that $\mathbb E \widetilde{Y}_{l_1}\widetilde{Y}_{l_2}\leq \epsilon$ whenever $l_1\neq l_2$. The heart of the matter lies in the following proposition. 
\begin{prop}\label{prop:CountAt4Scales}
    Assume that $H$ is as in Theorem \ref{thm:asBoundsShort}. Let $l_1$ and $l_2$ be  primes in the range $h(X)/2 < l_1,  l_2\leq h(X)=(X/H(X))^{\delta}$. The number of non-trivial $n_1, n_3\in\mathcal B_{N_{l_1}}$ and $n_2, n_4\in\mathcal B_{N_{l_2}}$ with $P^+(n_1)=P^+(n_2)$ and $P^+(n_3)=P^+(n_4)$ such that $n_1n_2n_3n_4=\square$ (or $n_1n_4=n_2n_3$) is $\ll H_{l_1}H_{l_2}/h(X)^2$ if $\delta>0$ and $\varepsilon>0$ are small enough.\footnote{Recall the definition of $\mathcal B_{N_l}$ for the dependence on $\varepsilon$.}
\end{prop}

\begin{proof}
Write $\gamma_l = H_l/N_l$ so that $H_l = \gamma_l N_l$. Notice that $$\gamma_{l_1}\asymp \gamma_{l_2}\asymp \gamma = \gamma(X)= H(Xh(X))/(Xh(X))$$ uniformly in the values of $l_1, l_2$, since $H$ is increasing and concave. Without loss of generality, we may assume that $l_1\leq  l_2$. 

We begin with the Steinhaus case. We utilize the parametization of the solutions to $n_1n_4=n_2n_3$ as 
\begin{equation*}
    n_1= ga, n_2 = gb, n_3 =ha, n_4 = hb,
\end{equation*} where $(a, b)=1$ and $g = (n_1, n_2), h=(n_3, n_4).$ In particular, $P^+(n_1)\mid g$ and $P^+(n_3) \mid h$, so both $g, h > (Xh(X))^{\frac{2}{3}}.$ 
Since we are concerned with non-trivial solutions, we may assume that $g\neq h$ and $a\neq b$.  
Note that 
\begin{equation*}
(1-\gamma_{l_1})^2 = \frac{l_2^2N_{l_1}^2(1-\gamma_{l_1})^2}{l_1^2N_{l_2}^2}\leq\frac{l_2^2n_1n_3}{l_1^2n_2n_4} = \frac{(l_2a)^2}{(l_1b)^2}\leq \frac{l_2^2N_{l_1}^2}{l_1^2N_{l_2}^2(1-\gamma_{l_2})^2} = \frac{1}{(1-\gamma_{l_2})^2},
\end{equation*}
thus 
\begin{equation}\label{eq:a,b,gamma}
    \gamma\ll \left\lvert\frac{l_2a}{l_1b} - 1\right\rvert\ll \gamma.
\end{equation} In particular, if $l_2a\neq l_1b$, then we have both $l_2a\gg \gamma^{-1}$ and $l_1b\gg \gamma^{-1}$. Since $l_1, l_2$ are distinct primes and $(a, b) = 1$, the only way we could have $l_2a=l_1b$ is if $a=l_1, b=l_2$. In this case, the number of choices for $g$ and $h$ is bounded by $H_{l_1}H_{l_2}/(l_1l_2)\ll H_{l_1}H_{l_2}/h(X)^2$. Throughout now we may assume that $l_2a\gg \gamma^{-1}$ and $l_1b\gg \gamma^{-1}$. 

Let us first handle the case when $H_1\leq N_1/(\log N_1)^2$. Given $a, b\ll N_{l_1}^{\frac{1}{3}}$, the number of choices for $g$ and $h$ satisfying the constraints for the $n_i$ is certainly bounded by $H_{l_1}H_{l_2}/(ab)$. Note that $b\gg l_2a/l_1$  by \eqref{eq:a,b,gamma}, therefore the number of solutions with $l_2a, l_1b\gg \gamma^{-1}$ is bounded by 
\begin{align*}
H_{l_1}H_{l_2}\sum_{\substack{a, b \ll N_{l_1}^{\frac{1}{3}} \\ l_2a, l_1b\gg \gamma^{-1} \\ \eqref{eq:a,b,gamma}} } \frac{1}{ab} & \ll H_{l_1}H_{l_2} \sum_{\substack{a\ll N_{l_1}^{\frac{1}{3}} \\ a\gg (l_2\gamma)^{-1}}} \frac{1}{a^2}\sum_{b : \eqref{eq:a,b,gamma}}1 
\\ &  \ll H_{l_1}H_{l_2}\sum_{\substack{a\ll N_{l_1}^{\frac{1}{3}} \\ a\gg (l_2\gamma)^{-1}}}  \frac{1}{a^2}\left(\gamma l_2a/l_1 +1 \right)    \ll \gamma H_{l_1}H_{l_2}(\log X + l_2) \\ & \ll \gamma H_{l_1} H_{l_2} h(X),
\end{align*}
since $h(X)\gg \log X$ in this case. Accounting for the solutions with $l_2a=l_1b$, we conclude that the number of non-trivial solutions is bounded by $H_{l_1}H_{l_2}/h(X)^2$. 

Let us now turn to the case when $H_1>N_1/(\log N_1)^2$. Let $K=(1+\varepsilon)\log\log X$, so that $\Omega(gahb)\leq 2K$. The number of non-trivial solutions to the equation $n_1n_4=n_2n_3$ with $a\neq l_1, b\neq l_2$ is then 
\begin{equation}\label{eq:SteinhausLambda}\ll2^{2K}\sum_{\substack{a\neq b \\  \gamma \ll \left\lvert\frac{l_2a}{l_1b}-1\right\rvert \ll \gamma \\ l_2a, l_1b \gg\gamma^{-1}\\ a, b\ll N_{l_1}^{\frac{1}{3}}}} 2^{-\Omega(ab)}\sum_{\substack{g\neq h \\ g, h\gg {N_{l_1}}^{\frac{2}{3}} \\ \max\{N_{l_1}(1-\gamma)/a, N_{l_2}(1-\gamma)/b\}\leq g, h\leq \min\{N_{l_1}/a, N_{l_2}/b \}}}2^{-\Omega(gh)}.
\end{equation}
Applying Proposition \ref{prop:nair-tenenbaum} to the sums over $g$ and $h$ yields that the right-hand side of \eqref{eq:SteinhausLambda} is
\begin{equation}
\label{eq:SteinhausLambda2}
 \ll 2^{2K} \frac{H_{l_1}H_{l_2}}{\log X}\sum_{\substack{a\neq b \\  \gamma \ll \left\lvert\frac{l_2a}{l_1b}-1\right\rvert \ll \gamma \\ l_2a, l_1b \gg\gamma^{-1}\\ a, b\ll N_{l_1}^{\frac{1}{3}}}} \frac{2^{-\Omega(ab)}}{ab}  \ll 2^{2K} \frac{H_{l_1}H_{l_2}}{\log X}\sum_{\substack{l_2a \gg\gamma^{-1}\\ a\ll N_{l_1}^{\frac{1}{3}}}} \frac{2^{-\Omega(a)}}{a^2} \sum_{\substack{a\neq b \\  \gamma \ll \left\lvert\frac{l_2a}{l_1b}-1\right\rvert \ll \gamma \\ l_1b \gg\gamma^{-1} }} 2^{-\Omega(b)}.
\end{equation}
If $l_2a\ll \gamma^{-2}$, we simply bound the sum over $b$ by $\gamma l_2a\gg 1$, namely the number of terms, and then the sum over such $a$ is 
\begin{equation*}
    \ll\gamma l_2 \sum_{l_2a\ll \gamma^{-2}}\frac{1}{a} \ll \gamma l_2 \log \gamma^{-1}.
\end{equation*}
If $l_2a\gg \gamma^{-2}$, we may apply Proposition \ref{prop:nair-tenenbaum} to the sum over $b$, yielding that the right-hand side of    \eqref{eq:SteinhausLambda2} is
\begin{equation}\label{eq:SteinhausLambda3}
  \ll 2^{2K} \gamma l_2\frac{H_{l_1}H_{l_2}}{(\log X)^{2\frac{1}{2}}} \left(\log\gamma^{-1} + \sum_{\substack{l_2a\gg \gamma^{-2} \\ a\ll N_{l_1}^{\frac{1}{3}}}} \frac{2^{-\Omega(a)}}{a(\log (l_2a))^{\frac{1}{2}}} \right).
\end{equation}
By a standard dyadic decomposition for $a$ and using Proposition \ref{prop:nair-tenenbaum} again, we obtain 
\begin{equation*}
    \sum_{a\ll N_{l_1}^{\frac{1}{3}}} \frac{2^{-\Omega(a)}}{a(\log a)^{\frac{1}{2}}} \ll \sum_{A=2^k\ll X^2 }\frac{1}{A (\log A)^{\frac{1}{2}}}\sum_{A<a\leq 2A }2^{-\Omega(a)}  \ll \sum_{k\ll \log X}\frac{1}{k}\ll \log \log X.
\end{equation*}
 Putting everything together (including the solutions with $a=l_1, b=l_2$), we conclude that the number of non-trivial solutions is 
\begin{equation*}
    \ll \frac{\gamma h(X)H_{l_1}H_{l_2}}{(\log X)^{1-2(1+\varepsilon)\log 2}}\log\log X + \frac{H_{l_1}H_{l_2}}{h(X)^2}\ll \frac{H_{l_1}H_{l_2}}{h(X)^2}
\end{equation*}
by making $\delta>0$ smaller if necessary, since $H(X)\leq X/(\log X)^c$ with $c>2\log 2 -1$. 


Next, we consider the Rademacher case, which is very similar but a bit more technical. We utilize the parametrization of the solutions to $n_1n_2n_3n_4=\square$ as 
    \begin{equation*}
        n_1 = Aru, n_2 = Asv, n_3 =Brv, n_4 = Bsu,
    \end{equation*}
    where the variables $A, B, r, s, u, v$ satisfy $A = (n_1, n_2), B = (n_3, n_4)$, and $r, u, s, v$ parametrize the solutions to $(n_1/A)(n_4/B) = (n_2/A)(n_3/B)$, as in the Steinhaus case. In particular, the variables $r, u, s, v$ are all pairwise coprime. We may assume that $P^+(n_1)\neq P^+(n_3)$ (the contribution of such $n_i$ is very small, see the bound above in (I')), so that $A\neq B$ and $A, B\gg N_{l_1}^{\frac{2}{3}}$. Note that 
    \begin{equation*}
        (1-\gamma_{l_1})^2\leq \frac{l_2^2n_1n_3}{l_1^2 n_2n_4} = \frac{(l_2r)^2}{(l_1s)^2}\leq \frac{1}{(1-\gamma_{l_2})^2},
    \end{equation*} thus 
    \begin{equation}\label{eq:r,s,gamma}
        \gamma \ll \left\lvert \frac{l_2r}{l_1s} - 1 \right\rvert \ll \gamma.
    \end{equation} In particular if $l_2r\neq l_1s$ then $l_2r, l_1s\gg \gamma^{-1}$. Furthermore, since $(r, s)=1$, the only way we could have $l_2r=l_1s$ is if $r=l_1, s=l_2$. In this case, note that $m_1 = n_1/r, m_3 =n_3/r, m_2=n_2/s, m_4=n_2/s$ satisfy $m_1m_3=m_2m_4$ and all of the $m_i$ lie in an interval of length $\ll \gamma X$ centered at $X$ (since $N_{l_1}/r = X=N_{l_2}/s$). However, the number of choices for such $m_i$ is $\ll (\gamma X)^2\ll H_{l_1}H_{l_2}/(rs)\ll H_{l_1}H_{l_2}/h(X)^2$ by our work above in the Steinhaus case, or the work of Soundararajan and Xu \cite{soundararajan2023clt}. Henceforth we may assume that $l_2r\neq l_1s$, thus $l_2r, l_1s\gg \gamma^{-1}$. 
    
    Moreover
    \begin{equation*}
        (1-\gamma_{l_1})(1-\gamma_{l_2})\leq \frac{n_1n_4}{n_2n_3} = \frac{u^2}{v^2}\leq \frac{1}{(1-\gamma_{l_1})(1-\gamma_{l_2})},
    \end{equation*} thus 
    \begin{equation}
        \label{eq:u, v, gamma}
        \gamma \ll \left\lvert \frac{u}{v}-1\right\rvert\ll \gamma.
    \end{equation}
    Again, if $u\neq v$, we have $u, v\gg \gamma^{-1}$. The cases where $r=s$ or $u=v$ (which imply $r=s=1$ or $u=v=1$, respectively, due to coprimality) simply reduce to our earlier work in the Steinhaus case. Therefore, we may now assume that $u\neq v, r\neq s, l_2r\neq l_1s$ and $u, v, l_2r, l_1s\gg\gamma^{-1}$. The number of such solutions is bounded by 
    \begin{equation}\label{eq:RademacherLambda1}
        2^{2K}\sum_{\substack{r, u, s, v \\ l_2r, l_1s, u, v\gg \gamma^{-1} \\ ru, rs, su, sv\leq N_{l_2}^{\frac{1}{3}} \\ \eqref{eq:r,s,gamma}, \eqref{eq:u, v, gamma}}} 2^{-\Omega(rusv)}\sum_{\substack{A, B\\ A, B\gg N_{l_2}^{\frac{1}{3}} \\ (N_{l_1} - H_{l_1})/(ru) \leq A\leq  N_{l_1}/(ru) \\ (N_{l_2} - H_{l_2})/(su)\leq A\leq  N_{l_2}/(su)}} 2^{-\Omega(AB)}.
    \end{equation}
Proposition \ref{prop:nair-tenenbaum} applied  to the sums over $A$ and $B$ yields that the right-hand side of  \eqref{eq:RademacherLambda1} is
     \begin{equation}
         \label{eq:RademacherLambda2}
         \ll 2^{2K}\frac{H_{l_1}H_{l_2}}{\log X} \sum_{\substack{r, u, s, v \\ l_2r, l_1s, u, v\gg \gamma^{-1} \\ ru, rs, su, sv\leq N_{l_2}^{\frac{1}{3}} \\ \eqref{eq:r,s,gamma}, \eqref{eq:u, v, gamma}}} \frac{2^{-\Omega(rusv)}}{rsu^2}.
     \end{equation}
     Using \eqref{eq:r,s,gamma} we bound $s\gg l_2r/l_1\asymp r$, hence the right-hand side of       \eqref{eq:RademacherLambda2} is 
\begin{equation}\label{eq:RademacherLambda3}
    \ll 2^{2K}\frac{H_{l_1}H_{l_2}}{\log X} \sum_{\substack{r, u, s, v \\ l_2r, l_1s, u, v\gg \gamma^{-1} \\ ru, rs, su, sv\leq N_{l_2}^{\frac{1}{3}} \\ \eqref{eq:r,s,gamma}, \eqref{eq:u, v, gamma}}} \frac{2^{-\Omega(rusv)}}{r^2u^2}.
     \end{equation}
     Dropping all relations between the pairs $(r, s)$ and $(u, v)$, we observe that the bound in \eqref{eq:RademacherLambda3} is identical to the square of the sum over $a$ and $b$ in \eqref{eq:SteinhausLambda2}, thus may be bounded by $(\gamma l_2\log\log X)^2$. Accounting for the solutions with $l_2r=l_1s$ or $r=s=1$ or $u=v=1$, in a similar fashion as before, we conclude that the number of non-trivial solutions is 
     \begin{equation}
         \ll \frac{H_{l_1}H_{l_2}}{h(X)^2}
     \end{equation} by making $\delta>0$ and $\varepsilon>0$ smaller if necessary. 
    \end{proof}

 Applying Markov's inequality with the second moment and using Proposition \ref{prop:CountAt4Scales} gives
    \begin{equation*}
        \mathbb P\left(\mathbb E\widetilde{Y}_{l_1}\widetilde{Y}_{l_2}\geq \epsilon \right)\ll \frac{h(X)^{-2}}{\epsilon^2} \hspace{0.2cm}\text{ and }\hspace{0.2cm} \mathbb P\left(\left\lvert \mathbb EY_l^2-1\right\rvert\geq\epsilon \right)\ll \frac{h(X)^{-2}}{\epsilon^2}
    \end{equation*}
    for any $l_1\neq l_2$ and $l$. 
    For fixed and small enough $\epsilon >0$, we conclude that for fixed $l_1\neq l_2$, we have $\mathbb E \widetilde{Y}_{l_1}\widetilde{Y}_{l_2}\leq \epsilon$ and $\mathbb EY_l^2\geq 1-\epsilon$ with probability $1-O(h(X)^{-2}).$ Taking a union bound over $\asymp k^2$ values of $l_1\neq l_2$ and $l$ gives that $\mathbb E\widetilde{Y}_{l_1}\widetilde{Y}_{l_2}\leq \epsilon$ and $\mathbb EY_l^2\gg 1$ with probability $1- O(k^2h(X)^{-2})$. Now Lemma \ref{lemma:normalsBigLemma} gives 
    \[
\mathbb{P}\left(\max_{l} Y_l\ll \sqrt{\log k}\right) \leq  \exp\left(-\Theta\left(\frac{k^{O(1)}}{\sqrt{\log k}}\right)\right) + k^{-O(1)}
\]
for a small enough implied constant. Combining this with \eqref{eq:shortIntsNormal} with fixed $\eta>0$ gives that
\begin{equation*}\label{eq:shortIntsNormal2}
\max_l \widetilde{S}_{N_l} \gg \sqrt{\log k}
\end{equation*}
holds with probability 
\begin{equation*}
    1- O\left(k^{\frac{3}{2}}\eta^{-1} \left(k\sqrt{\varepsilon_1' + \varepsilon_2}\right)^{\frac1{5(k+1)}}  + \exp\left(-\Theta\left(\frac{k^{O(1)}}{\sqrt{\log k}}\right)\right) + k^{-O(1)}\right).
\end{equation*}
Recalling Lemmas \ref{lemma:ignoreTooManyPrimesShort} and \ref{lemma:ignoresmoothShort}, \eqref{eq:splitting1}, \eqref{eq:splitting2}, the values of $\varepsilon_1', \varepsilon_2$, and that $k = (X/H(X))^{\varepsilon_0}$ for a suitably small $\varepsilon_0>0$, we obtain 
\begin{equation*}
    \max_ l \left\lvert\frac{1}{\sqrt{H_l}}\sum_{N_l - H_l<n\leq N_l } f(n)\right\rvert \gg \sqrt{\log \frac{X}{H(X)}}
\end{equation*}
with probability $1- O((\log\log X)^{-\frac{2}{3}})$. Here we replaced the normalization by $\lvert\mathcal B_{N_l}\rvert$ in $\widetilde S_{N_l}$ with $\sqrt{H_l}$, as we may for $H\geq N^{\frac{11}{15}}$, in light of Lemmas \ref{lemma:tooManyPrimes} and \ref{lemma:unSmoothSquarefree}. This finishes the proof of Theorem \ref{thm:localizedBoundShort}.







%    Bibliographies can be prepared with BibTeX using amsplain,
%    amsalpha, or (for "historical" overviews) natbib style.
\bibliographystyle{amsplain}
\bibliography{ref.bib}
%    Insert the bibliography data here.

\end{document}

%%%%%%%%%%%%%%%%%%%%%%%%%%%%%%%%%%%%%%%%%%%%%%%%%%%%%%%%%%%%%%%%%%%%%%%%

%    Templates for common elements of a journal article; for additional
%    information, see the AMS-LaTeX instructions manual, instr-l.pdf,
%    included in the TRAN author package, and the amsthm user's guide,
%    linked from http://www.ams.org/tex/amslatex.html .

%    Section headings
\section{}
\subsection{}

%    Ordinary theorem and proof
\begin{theorem}[Optional addition to theorem head]
% text of theorem
\end{theorem}

\begin{proof}[Optional replacement proof heading]
% text of proof
\end{proof}

%    Figure insertion; default placement is top; if the figure occupies
%    more than 75% of a page, the [p] option should be specified.
\begin{figure}
\includegraphics{filename}
\caption{text of caption}
\label{}
\end{figure}

%    Mathematical displays; for additional information, see the amsmath
%    user's guide, linked from http://www.ams.org/tex/amslatex.html .

% Numbered equation
\begin{equation}
\end{equation}

% Unnumbered equation
\begin{equation*}
\end{equation*}

% Aligned equations
\begin{align}
  &  \\
  &
\end{align}

%-----------------------------------------------------------------------
% End of tran-l-template.tex
%-----------------------------------------------------------------------
